%!TEX root = ../../main.tex
% ============================================================
% 几何作图 / 分数涂色与等值分数
% 本文件由 main.tex 通过 \input 引入，请勿单独编译
% 重构日期: 2026-04-11
% ============================================================

\subsection{解答：在图中涂色表示分数}

\vspace{0.4cm}

\noindent\textbf{1. 在图中涂色表示它下面的分数。}

\vspace{0.8cm}

\begin{center}
\begin{tikzpicture}[>=Stealth, scale=0.85, line width=1pt] % 整体思路：三幅图分别通过“平均分后涂色”的方式表示分数，因此每一幅都先画外形或网格，再填充对应份数，最后在下方标分数

  % 1. 正方形网格 3x3
  \begin{scope}[shift={(0,0)}] % 第一幅图保持在左侧原位
    \fill[black!70] (0, 2) rectangle (3, 3); % 先把最上面一整行 3 个小正方形涂黑；rectangle 表示用对角点画矩形
    \fill[black!70] (0, 1) rectangle (2, 2); % 再把中间一行左边 2 个小格涂黑，这样一共涂了 5 个小格
    \draw[line width=0.8pt, step=1cm] (0,0) grid (3,3); % 画 3×3 网格；step=1cm 表示每一格边长 1cm；line width=0.8pt 为细边框
    \node[below=15pt, font=\Large] at (1.5, 0) {$\frac{5}{9}$}; % 在图下方 15pt 处写出分数 5/9
  \end{scope} % 结束第一幅正方形网格图

  % 2. 平行四边形
  \begin{scope}[shift={(5.5,0)}] % 第二幅图整体右移 5.5cm
    % 涂色区域 (左上角子平行四边形)
    \fill[black!70] (0.75, 1.5) -- (1.95, 1.5) -- (2.7, 3) -- (1.5, 3) -- cycle; % 用四个顶点围成一个小平行四边形并填黑，表示 4 份中的 1 份；cycle 自动闭合图形
    % 外框
    \draw[line width=0.8pt] (0,0) -- (2.4,0) -- (3.9,3) -- (1.5,3) -- cycle; % 画大平行四边形外框
    % 分割线
    \draw[line width=0.8pt] (1.2,0) -- (2.7,3); % 画一条与侧边平行的分割线，把图形纵向分成两份
    \draw[line width=0.8pt] (0.75,1.5) -- (3.15,1.5); % 再画一条水平样式的分割线，把图形横向分成两份，最终得到 4 等份
    % 分数
    \node[below=15pt, font=\Large] at (1.95, 0) {$\frac{1}{4}$}; % 在图下方标出对应分数 1/4
  \end{scope} % 结束第二幅平行四边形分数图

  % 3. 椭圆与圆圈
  \begin{scope}[shift={(13.5, 1.5)}] % 第三幅图右移并略上移，让椭圆与前两图对齐
    % 椭圆外框
    \draw[line width=0.8pt] (0, 0) ellipse (3.6cm and 1.6cm); % 以原点为中心画椭圆；横半轴 3.6cm，纵半轴 1.6cm，表示整体“1”的容器
    
    % 排列圆圈 (2行 x 5列)
    \foreach \x in {-2, -1, 0, 1, 2} { % 外层循环控制 5 列圆圈的位置
      \foreach \y in {-0.6, 0.6} { % 内层循环控制上下两行圆圈的位置
        \ifnum \x<1 % 条件判断：当前列编号若小于 1，就属于左边 3 列，需要涂黑
          % 涂黑的圆圈 (左边 3 列共 6 个)
          \filldraw[fill=black!85, draw=black, line width=0.8pt] (\x*1.1, \y) circle (0.35cm); % filldraw 同时画轮廓并填充；圆心坐标由 \x*1.1 和 \y 给出；circle (0.35cm) 表示半径 0.35cm
        \else
          % 空白圆圈 (右边 2 列共 4 个)
          \draw[line width=0.8pt] (\x*1.1, \y) circle (0.35cm); % 对右边两列只画空心圆，不填色
        \fi
      }
    }
    % 分数
    \node[below=15pt, font=\Large] at (0, -1.5) {$\frac{3}{5}$}; % 在椭圆下方写出分数 3/5
  \end{scope} % 结束第三幅“圆圈分组”分数图

\end{tikzpicture}
\end{center}

\vspace{0.6cm}
{\small\color{gray}
\textbf{解析：}\\
1. $\frac{5}{9}$：把正方形平均分成 $9$ 份，涂色部分占其中的 $5$ 份（涂满上方的 $3$ 个和中间行左侧的 $2$ 个小正方形）。\\
2. $\frac{1}{4}$：把平行四边形平均分成 $4$ 份，涂色部分占其中的 $1$ 份（本题涂满左上方的 $1$ 个平行四边形）。\\
3. $\frac{3}{5}$：共有 $10$ 个圆圈，平均分作 $5$ 份，每份 $2$ 个。$\frac{3}{5}$ 表示取其中的 $3$ 份，即 $2 \times 3 = 6$ 个圆圈，故将左侧的 $6$ 个小圆圈全部涂成黑色实体。
}

\vspace{0.8cm}


\vspace{0.6cm}

{\small\color{blue!70!black}
\textbf{绘图基本原理与步骤：}\\
\textbf{1. 坐标定位与几何构建}：根据题意中的已知长度和角度，使用 \texttt{coordinate} 定义关键顶点坐标，通过 \texttt{draw} 指令按序连接成完整几何图形。线宽、颜色参数区分原图（黑色）与答案（红色 \texttt{red!70!black}）图层。\\
\textbf{2. 辅助量标注}：利用 \texttt{node} 的方向锚点（\texttt{above}、\texttt{below}、\texttt{left}、\texttt{right}）将角度、边长、面积等数据标注在图形周围，避免与线条或其他标注重叠。若涉及角弧则使用 \texttt{arc} 或 \texttt{pic[angle]} 命令渲染。\\
\textbf{3. 虚实线分层}：题干已知条件用实线绘制，解答新增的辅助线或操作痕迹用 \texttt{dashed} 虚线或彩色加粗线标识，确保读者一眼区分原有信息与作答内容。
}

%% ==============================
%% 分数、小数、百分数互化
%% ==============================
\newpage

\newpage

\subsection{解答：六边形的分数涂色}

\vspace{0.6cm}

\textbf{\LARGE 82.} \textbf{将每个六边形分别涂色表示下面的分数。\\[1ex]
第一图下标注：$\displaystyle\frac{1}{2}$ \qquad 第二图下标注：$\displaystyle\frac{1}{3}$ \qquad 第三图下标注：$\displaystyle\frac{1}{6}$}

\vspace{0.4cm}
\textbf{解：}

\textbf{分析：}
如图所示，每个正六边形都被内部相交的三条对角线平均分成了 $6$ 个完全相同的小正三角形。
要表示特定的分数，只需要计算应涂色的小三角形数量：
1. 第一图表示 $\displaystyle\frac{1}{2}$：需要涂色总块数的 $\displaystyle\frac{1}{2}$，即 $6 \times \frac{1}{2} = 3$ 块。
2. 第二图表示 $\displaystyle\frac{1}{3}$：需要涂色总块数的 $\displaystyle\frac{1}{3}$，即 $6 \times \frac{1}{3} = 2$ 块。
3. 第三图表示 $\displaystyle\frac{1}{6}$：需要涂色总块数的 $\displaystyle\frac{1}{6}$，即 $6 \times \frac{1}{6} = 1$ 块。

\vspace{0.6cm}
\begin{center}
\begin{tikzpicture}[>=Stealth, scale=1.6]

  % 定义绘制一个被6等分的正六边形的宏
  % 参数1: x位移
  \def\drawHexagon#1{
    \begin{scope}[shift={(#1, 0)}]
      % 画六边形外框
      \draw[line width=1pt] (0:1) \foreach \i in {1,2,3,4,5} { -- (\i*60:1) } -- cycle;
      % 画三条虚线对角线
      \draw[dashed, line width=0.6pt] (0:1) -- (180:1);
      \draw[dashed, line width=0.6pt] (60:1) -- (240:1);
      \draw[dashed, line width=0.6pt] (120:1) -- (300:1);
    \end{scope}
  }

  % 第一个图形：涂 3 块 (1/2)
  \begin{scope}[shift={(0,0)}]
    % 填充 0~180 度的 3 个小正三角形
    \fill[blue!30] (0,0) -- (0:1) \foreach \i in {1,2,3} { -- (\i*60:1) } -- cycle;
    \drawHexagon{0}
    \node[below, font=\Large] at (0, -1.2) {$\displaystyle\frac{1}{2}$};
  \end{scope}

  % 第二个图形：涂 2 块 (1/3)
  \begin{scope}[shift={(3.5,0)}]
    % 填充 60~180 度的 2 个小正三角形 (顶部及左上)
    \fill[blue!30] (0,0) -- (60:1) \foreach \i in {2,3} { -- (\i*60:1) } -- cycle;
    \drawHexagon{0}
    \node[below, font=\Large] at (0, -1.2) {$\displaystyle\frac{1}{3}$};
  \end{scope}

  % 第三个图形：涂 1 块 (1/6)
  \begin{scope}[shift={(7,0)}]
    % 填充 60~120 度的 1 个小正三角形 (正上方)
    \fill[blue!30] (0,0) -- (60:1) -- (120:1) -- cycle;
    \drawHexagon{0}
    \node[below, font=\Large] at (0, -1.2) {$\displaystyle\frac{1}{6}$};
  \end{scope}

\end{tikzpicture}
\end{center}

\vspace{0.4cm}
{\small\color{blue!70!black}
\textbf{绘图原理与步骤：}\\
\textbf{1. 极坐标正多边形生成：} 利用极坐标系循环生成法 \texttt{(A:R)} 快速渲染完全标准正六边形，设定外接圆定长半径为 $1$。并人为规定顶点的极角依次为 $0^\circ, 60^\circ, 120^\circ \dots 300^\circ$。这一设定不仅确保了图形的顶边与底边水平，且使得连接 $0^\circ$ 和 $180^\circ$ 的对角线刚好是一根的水平中轴线，与原题素材的物理放置姿态完全锲合。\\
\textbf{2. 均等基本块划分：} 通过追加 \texttt{\textbackslash draw[dashed]} 指令连接全部核心对角线，实现在渲染外框的同时，在图形内部打通了隔断线。物理上将原本致密的整块大正六边形成功均分为 $6$ 块等面积的正三角形阵列单元。\\
\textbf{3. 分割填色法则：} 将题目提出的代数分数逻辑映射为涂色数量逻辑（$1/2\to 3$ 块，$1/3\to 2$ 块，$1/6\to 1$ 块）。然后启用特定的区块填充层 \texttt{\textbackslash fill[blue!30]}，精确约束中心 $(0,0)$ 延伸到外围对应节点之间的扇形发散式多边形面域，彻底杜绝手工填充导致的色彩溢出或留白死角。
}

%% ==============================
%% 第92题（补充题）：用不同方法表示正方形的 1/4
%% ==============================
\newpage

\newpage

\subsection{解答：表示分数的多种折叠模型}

\vspace{0.6cm}

\textbf{\LARGE 83.} \textbf{（第 4 题）你能用不同的方法表示每个正方形的 $\displaystyle\frac{1}{4}$ 吗？先用纸折一折，再画一画。}

\vspace{0.4cm}
\textbf{解：}

\textbf{分析：}
要把正方形平均分成 $4$ 份并着色表示出其中的 $\displaystyle\frac{1}{4}$，常见的基于对称性的折法（等分法）共有以下 4 种典型图案：
1. \textbf{竖向等分}：将正方形往同一方向连续对折两次，得到 4 个相同高度的竖直长条。涂色其中 1 块。
2. \textbf{横向等分}：与竖向等分同理，得到 4 个相同宽度的横向长条。涂色其中 1 块。
3. \textbf{十字等分}：横向和纵向各中心对折一次，得到 4 个相同的小正方形。涂色其中 1 块小正方形。
4. \textbf{对角线等分}：沿着正方形的两条对角线分别对折，得到 4 个相同的等腰直角三角形。涂色其中 1 块。

\vspace{0.6cm}
\begin{center}
\begin{tikzpicture}[>=Stealth, x=2.2cm, y=2.2cm]

  % 提取绘制单体正方形的宏
  \def\drawSquare#1{
    % 画实线外框
    \draw[line width=1pt] (#1, 0) rectangle (#1+1, 1);
  }

  % 1. 竖向四等分
  \begin{scope}[shift={(0,0)}]
    % 涂色第一条
    \fill[blue!20] (0, 0) rectangle (0.25, 1);
    \drawSquare{0}
    % 画分划虚线
    \draw[dashed, line width=0.6pt] (0.25, 0) -- (0.25, 1);
    \draw[dashed, line width=0.6pt] (0.5,  0) -- (0.5,  1);
    \draw[dashed, line width=0.6pt] (0.75, 0) -- (0.75, 1);
    \node[below, font=\small] at (0.5, -0.25) {方法一};
  \end{scope}

  % 2. 纵向四等分（横向长条）
  \begin{scope}[shift={(1.4,0)}]
    % 涂色最下一条
    \fill[blue!20] (0, 0) rectangle (1, 0.25);
    \drawSquare{0}
    \draw[dashed, line width=0.6pt] (0, 0.25) -- (1, 0.25);
    \draw[dashed, line width=0.6pt] (0, 0.5)  -- (1, 0.5);
    \draw[dashed, line width=0.6pt] (0, 0.75) -- (1, 0.75);
    \node[below, font=\small] at (0.5, -0.25) {方法二};
  \end{scope}

  % 3. 十字四等分
  \begin{scope}[shift={(2.8,0)}]
    % 涂色左下角小正方形
    \fill[blue!20] (0, 0) rectangle (0.5, 0.5);
    \drawSquare{0}
    \draw[dashed, line width=0.6pt] (0.5, 0) -- (0.5, 1);
    \draw[dashed, line width=0.6pt] (0, 0.5) -- (1, 0.5);
    \node[below, font=\small] at (0.5, -0.25) {方法三};
  \end{scope}

  % 4. 对角线四等分
  \begin{scope}[shift={(4.2,0)}]
    % 涂色下方三角形
    \fill[blue!20] (0, 0) -- (1, 0) -- (0.5, 0.5) -- cycle;
    \drawSquare{0}
    \draw[dashed, line width=0.6pt] (0, 0) -- (1, 1);
    \draw[dashed, line width=0.6pt] (0, 1) -- (1, 0);
    \node[below, font=\small] at (0.5, -0.25) {方法四};
  \end{scope}

\end{tikzpicture}
\end{center}

\vspace{0.4cm}
{\small\color{blue!70!black}
\textbf{绘图原理与步骤：}\\
\textbf{1. 图元定位阵列：} 建立坐标域，定义基础正方形内部比例常数为 $1$。分别于 $x=0, 1.4, 2.8, 4.2$ 的偏移处并行铺设四个基准大实体正方形外框，还原原题干中横向并排依次排开的四大方框环境。\\
\textbf{2. 几何代数等分模型：} 展现物理操作的“折纸”效果即是向框架内追加平分级的辅助线。针对等分这一严密的数学命题，分别编写了三等距垂直律、水平分割律、正交十字对偶中心律以及点对称极限折痕（交叉对角）等四种拓扑方程，通过 \texttt{dashed} 属性赋予拆分折痕。\\
\textbf{3. 独立集块高亮着色：} 利用 \texttt{\textbackslash fill[blue!20]} 的特定区域填充功能，对于矩形切块严格采用 \texttt{rectangle} 参数矩阵锁死，对于三角切块则采用闭环多边形连线 \texttt{cycle} 限位，强行高亮出四种形态截然不同但面积恒等占比 $1/4$ 的比例单元。
}


%% ==============================
%% 第93题（补充题）：分数涂色与大小比较
%% ==============================
\vspace{1.5cm}

\textbf{\LARGE 84.} \textbf{（第 3 题）先涂色表示分数，再比较每组分数的大小。}

\vspace{0.4cm}
\textbf{解：}

\textbf{分析：}
比较给定的两组真分数：
1. 第一组：同分母分数 $\displaystyle\frac{3}{8}$ 与 $\displaystyle\frac{5}{8}$ 比较。分母相同代表切分的标准块大小一样，此时“分子越大，分数实际所占的面积就越大”，故 $\displaystyle\frac{3}{8} < \displaystyle\frac{5}{8}$。
对应在图上的视觉直观操作：每个大正方形已被细分为 $8$ 个等分小三角形，在左图中涂去 $3$ 个，在右图中涂去 $5$ 个，清晰可见右侧涂色区更大。
2. 第二组：同分母分数 $\displaystyle\frac{3}{4}$ 与 $\displaystyle\frac{2}{4}$ 比较。同理，由分子数值可判明 $\displaystyle\frac{3}{4} > \displaystyle\frac{2}{4}$。
在图上的动作：每个大正方形已被切分为 $4$ 个等分小方格，左图按分子占涂 $3$ 个，右图涂 $2$ 个即见大小高低。

\vspace{0.6cm}
\begin{center}
\begin{tikzpicture}[>=Stealth, x=1.8cm, y=1.8cm]

  %% 第一组：3/8 与 5/8
  \begin{scope}[shift={(0,0)}]
    % 左边正方形 (3/8)
    \begin{scope}[shift={(0,0)}]
      % 涂 3 个小三角形
      \fill[blue!30] (0.5, 0.5) -- (0,0) -- (0.5,0) -- cycle;
      \fill[blue!30] (0.5, 0.5) -- (0.5,0) -- (1,0) -- cycle;
      \fill[blue!30] (0.5, 0.5) -- (1,0) -- (1,0.5) -- cycle;
      
      % 外框
      \draw[line width=1pt] (0,0) rectangle (1,1);
      % 米字分割线
      \draw[dashed, line width=0.6pt] (0.5,0) -- (0.5,1);
      \draw[dashed, line width=0.6pt] (0,0.5) -- (1,0.5);
      \draw[dashed, line width=0.6pt] (0,0) -- (1,1);
      \draw[dashed, line width=0.6pt] (0,1) -- (1,0);
    \end{scope}

    % 右边正方形 (5/8)
    \begin{scope}[shift={(1.4,0)}]
      % 涂 5 个小三角形
      \fill[blue!30] (0.5, 0.5) -- (0,0) -- (0.5,0) -- cycle;
      \fill[blue!30] (0.5, 0.5) -- (0.5,0) -- (1,0) -- cycle;
      \fill[blue!30] (0.5, 0.5) -- (1,0) -- (1,0.5) -- cycle;
      \fill[blue!30] (0.5, 0.5) -- (1,0.5) -- (1,1) -- cycle;
      \fill[blue!30] (0.5, 0.5) -- (1,1) -- (0.5,1) -- cycle;
      
      % 外框
      \draw[line width=1pt] (0,0) rectangle (1,1);
      % 米字分割线
      \draw[dashed, line width=0.6pt] (0.5,0) -- (0.5,1);
      \draw[dashed, line width=0.6pt] (0,0.5) -- (1,0.5);
      \draw[dashed, line width=0.6pt] (0,0) -- (1,1);
      \draw[dashed, line width=0.6pt] (0,1) -- (1,0);
    \end{scope}

    % 下方分数及比较圆圈符号
    \node[anchor=mid, font=\Large] at (0.5, -0.6) {$\displaystyle\frac{3}{8}$};
    \draw[line width=0.8pt] (1.2, -0.6) circle (0.2);
    \node[anchor=mid, font=\large] at (1.2, -0.6) {$<$};
    \node[anchor=mid, font=\Large] at (1.9, -0.6) {$\displaystyle\frac{5}{8}$};
  \end{scope}


  %% 第二组：3/4 与 2/4
  \begin{scope}[shift={(4.5,0)}]
    % 左边正方形 (3/4)
    \begin{scope}[shift={(0,0)}]
      % 涂 3 个小正方形
      \fill[blue!30] (0,0) rectangle (1,0.5);
      \fill[blue!30] (0,0.5) rectangle (0.5,1);
      
      % 外框
      \draw[line width=1pt] (0,0) rectangle (1,1);
      % 十字分割线
      \draw[dashed, line width=0.6pt] (0.5,0) -- (0.5,1);
      \draw[dashed, line width=0.6pt] (0,0.5) -- (1,0.5);
    \end{scope}

    % 右边正方形 (2/4)
    \begin{scope}[shift={(1.4,0)}]
      % 涂 2 个小正方形
      \fill[blue!30] (0,0) rectangle (1,0.5);
      
      % 外框
      \draw[line width=1pt] (0,0) rectangle (1,1);
      % 十字分割线
      \draw[dashed, line width=0.6pt] (0.5,0) -- (0.5,1);
      \draw[dashed, line width=0.6pt] (0,0.5) -- (1,0.5);
    \end{scope}

    % 下方分数及比较圆圈符号
    \node[anchor=mid, font=\Large] at (0.5, -0.6) {$\displaystyle\frac{3}{4}$};
    \draw[line width=0.8pt] (1.2, -0.6) circle (0.2);
    \node[anchor=mid, font=\large] at (1.2, -0.6) {$>$};
    \node[anchor=mid, font=\Large] at (1.9, -0.6) {$\displaystyle\frac{2}{4}$};
  \end{scope}

\end{tikzpicture}
\end{center}

\vspace{0.4cm}
{\small\color{blue!70!black}
\textbf{绘图原理与步骤：}\\
\textbf{1. 域隔离对比系：} 利用两个独立包裹的 \texttt{scope} （平移量相隔足够宽距）宏观分区建立了两套独立的对照绘图区。左侧为八等分体系，右侧为基准十字四等分体系，保持所有方块外延尺寸严密一致并排布局，彻底实现原题左右互不干扰的并排比列版面。\\
\textbf{2. 细分网格裁切场：} 对于“涂色表示分数的实质计算”，必须对空白基板做等效切分：左侧区域通过对角线加横竖中位线簇（施加 \texttt{dashed} 点划线属性）切割成 $8$ 个“全边等底正三角”；右侧区域则仅走单纯的十字正中切分分离为 $4$ 个纯正方微块块，为底层的数据结构绑定图元。\\
\textbf{3. 分区面积映射法则：} 紧扣两组代数算子中的分子系数（$3, 5$ 及 $3, 2$），指令程序采用天蓝色（\texttt{blue!30}）精确检索并吞噬掉对应个数的微型切块。并在最后一行布置底注级渲染——用中心坐标 \texttt{anchor=mid} 精准压平实线评判空圆圈，同时往圆心中注入机器级准确判断的 $<$ 及 $>$ 等比例排版的矢量不等号。
}

%% ==============================
%% 第94题（补充题）：等值分数涂色
%% ==============================
\newpage

\newpage

\subsection{解答：等值分数的涂色与转化}

\vspace{0.6cm}

\textbf{\LARGE 85.} \textbf{先在图中涂色表示已知分数的等值分数，再填空。}

\vspace{0.4cm}
\textbf{解：}

\textbf{分析：}
“等值分数”是指虽然分子分母的数字不同，但其表示的实质数值（在一个固定大小的图形中即指代对应相等的覆盖面积）永远相等的分数。
1. 第一个图形组：左侧的六边形被极点射线分成了 $3$ 个并排的菱形，阴影部分占据 $1$ 个菱形，即表示 $\displaystyle\frac{1}{3}$。为了保证面积的“等值”，右侧的六边形（已被全等切分为 $6$ 个小正三角形）其涂色区域必须与左图一致（占据左半边），此时刚好吞并了 $2$ 块小三角形，故得出等价分数公式 $\displaystyle\frac{1}{3} = \frac{2}{6}$。
2. 第二个图形组：左侧的正方形横向等分切成了 $4$ 份水平层，顶层被涂成了灰色，客观表示 $\displaystyle\frac{1}{4}$。而在右侧的正方形中，整个画面被进一步碎化成 $4 \times 4 = 16$ 个同尺寸的小方格，要想使得涂色面积与左边高度吻合一致，那就必须涂满第一排平行的 $4$ 个小方格，由此得出 $\displaystyle\frac{1}{4} = \frac{4}{16}$。

\vspace{0.6cm}
\begin{center}
\begin{tikzpicture}[>=Stealth, line width=0.8pt]

  %% ========== 1. 左侧六边形组 (1/3 和 2/6) ==========
  \begin{scope}[shift={(0,0)}]
    \def\R{1.4}
    % 左一：涂灰色 (题干要求)
    \fill[gray!30] (0,0) -- (120:\R) -- (180:\R) -- (240:\R) -- cycle;
    % 边界底侧六边形
    \draw (0:\R) \foreach \i in {1,...,5} { -- (\i*60:\R) } -- cycle;
    % 内部三连平分虚线
    \draw[dashed, line width=0.6pt] (0,0) -- (0:\R);
    \draw[dashed, line width=0.6pt] (0,0) -- (120:\R);
    \draw[dashed, line width=0.6pt] (0,0) -- (240:\R);
  \end{scope}

  % 左侧等式推导： 1/3 = 2/6
  \node[anchor=mid, font=\Large] at (1.8, -2.1) {$\displaystyle\frac{1}{3} \quad = \quad \frac{(\quad \textcolor{blue!80!black}{2} \quad)}{(\quad \textcolor{blue!80!black}{6} \quad)}$};

  \begin{scope}[shift={(3.6,0)}]
    \def\R{1.4}
    % 左二：涂蓝色作为解答图
    \fill[blue!30] (0,0) -- (120:\R) -- (180:\R) -- (240:\R) -- cycle;
    \draw (0:\R) \foreach \i in {1,...,5} { -- (\i*60:\R) } -- cycle;
    % 对角交叉虚线 (六等分)
    \draw[dashed, line width=0.6pt] (0:\R) -- (180:\R);
    \draw[dashed, line width=0.6pt] (60:\R) -- (240:\R);
    \draw[dashed, line width=0.6pt] (120:\R) -- (300:\R);
  \end{scope}


  %% ========== 2. 右侧正方形组 (1/4 和 4/16) ==========
  \begin{scope}[shift={(7.5,-1.2)}]
    \def\S{2.4}
    % 右一：涂灰色 (最顶层1/4)
    \fill[gray!30] (0, 0.75*\S) rectangle (\S, \S);
    \draw (0,0) rectangle (\S, \S);
    \draw[dashed, line width=0.6pt] (0, 0.25*\S) -- (\S, 0.25*\S);
    \draw[dashed, line width=0.6pt] (0, 0.5*\S)  -- (\S, 0.5*\S);
    \draw[dashed, line width=0.6pt] (0, 0.75*\S) -- (\S, 0.75*\S);
  \end{scope}

  % 右侧等式推导： 1/4 = 4/16
  \node[anchor=mid, font=\Large] at (10.6, -2.1) {$\displaystyle\frac{1}{4} \quad = \quad \frac{(\quad \textcolor{blue!80!black}{4} \quad)}{(\quad \textcolor{blue!80!black}{16} \quad)}$};

  \begin{scope}[shift={(11.3,-1.2)}]
    \def\S{2.4}
    % 右二：涂天蓝色解答 (最上面一排共计 4 个小格子)
    \fill[blue!30] (0, 0.75*\S) rectangle (\S, \S);
    \draw (0,0) rectangle (\S, \S);
    % 内部密织 16 宫格虚线阵列
    \foreach \i in {0.25, 0.5, 0.75} {
      \draw[dashed, line width=0.6pt] (0, \i*\S) -- (\S, \i*\S);
      \draw[dashed, line width=0.6pt] (\i*\S, 0) -- (\i*\S, \S);
    }
  \end{scope}

\end{tikzpicture}
\end{center}

\vspace{0.4cm}
{\small\color{blue!70!black}
\textbf{绘图原理与步骤：}\\
\textbf{1. 极坐标正方多边形衍生：} 沿袭运用的极坐标生成宏构建底侧的基础正六边形，结合虚线平分技术实现 $1/3$ 到 $1/6$ 的精密切割过渡；在此环节中，原题提供基准图被统一施配灰色 \texttt{gray!30} 透明图层充当环境前提背景。\\
\textbf{2. 欧几里得网格切分场：} 根据正方形切片映射原理，我们对相同基底尺寸（宽幅 $2.4\text{ cm}$）的正方形采取了两套割裂动作逻辑。一是单纯运用 \texttt{y} 轴刻度线实施粗颗粒切划；二是使用了二维复套循环 \texttt{\textbackslash foreach} 将其进行极限十六级粉碎化制备。\\
\textbf{3. 核心解题反馈印记：} 在待作答的图形轮廓上严格用“主观解答色——淡蓝色 \texttt{blue!30}”追溯描摹出了完全等比的等积投影版块（譬如那的 $4$ 枚连排小方格）；与此同时借由底部的等式渲染结构，自然地将 $2/6$ 及 $4/16$ 利用代码着色一并填入圆润的小括号之中。
}






\newpage

\begin{center}
  {\large\bfseries 解答：以下面的射线为一条边，画出各个角}
\end{center}

\vspace{0.5cm}

\noindent
\begin{tikzpicture}[scale=1.0]

  %% ==============================
  %% 第一图：45°角（给定射线水平向右）
  %% ==============================
  \begin{scope}[xshift=0cm]
    % 顶点
    \coordinate (A) at (0,0);
    % 给定射线（水平向右）端点
    \coordinate (B) at (3,0);
    % 答案射线（与水平线成45°）端点
    \coordinate (C) at ({3*cos(45)},{3*sin(45)});

    % 绘制给定射线（有端点，无箭头表示已给定的边）
    \draw[given ray] (A) -- (B);
    % 绘制答案射线（从顶点出发，45°方向）
    \draw[ray] (A) -- (C);

    % 顶点蓝点
    \node[vertex dot] at (A) {};

    % 角弧线
    \pic[draw=red!70!black, line width=0.8pt, angle radius=0.8cm,
         angle eccentricity=1.5, "$45°$", font=\small, color=red!70!black]
        {angle = B--A--C};

    % 标题
    \node[above] at (1.5, 2.4) {$45°$};
  \end{scope}

  %% ==============================
  %% 第二图：105°角（给定射线竖直向上）
  %% ==============================
  \begin{scope}[xshift=5.5cm]
    \coordinate (A) at (0,0);
    % 给定射线（竖直向上）
    \coordinate (B) at (0,3);
    % 答案射线方向：从竖直射线（90°）顺时针105°
    % 竖直向上=90°，顺时针105° => 90° - 105° = -15° (即从x轴正方向量取-15°)
    % 等效于从x轴正方向的 90°+105° = 195°（逆时针从右水平开始）
    % 实际上：顶角在竖直射线左侧，答案射线在右侧与给定射线夹角105°
    % 给定射线90°，若顺时针转105°，答案射线在90°-105°=-15°方向（即x轴负方向偏下15°）
    % 但通常画角是在凸侧：给定射线90°，答案射线在90°+105°=195° 或 90°-105°=-15°
    % 105°钝角，答案射线应在给定垂直射线的右侧，过渡到水平以下
    % 从给定射线(90°)顺时针转105°: 90°-105° = -15° => 即 345° = 方向(-15°)
    \coordinate (C) at ({3*cos(-15)},{3*sin(-15)});

    \draw[given ray] (A) -- (B);
    \draw[ray] (A) -- (C);
    \node[vertex dot] at (A) {};

    \pic[draw=red!70!black, line width=0.8pt, angle radius=0.7cm,
         angle eccentricity=1.6, "$105°$", font=\small, color=red!70!black]
        {angle = C--A--B};

    \node[above] at (0, 3.3) {$105°$};
  \end{scope}

  %% ==============================
  %% 第三图：165°角（比平角小15°，给定射线斜向右上方约45°）
  %% ==============================
  \begin{scope}[xshift=11cm]
    \coordinate (A) at (0,0);
    % 给定射线：斜向右上方（约45°）
    \coordinate (B) at ({2.5*cos(45)},{2.5*sin(45)});
    % 答案射线：与给定射线夹角165°
    % 给定射线方向45°，答案射线方向：45° + 165° = 210° 或 45° - 165° = -120°(=240°)
    % 取 45° - 165° = -120°，即在给定射线右侧顺时针165°
    % 从图示看给定射线斜右上，165°的另一边应在给定射线的左下侧附近
    \coordinate (C) at ({2.5*cos(45+165)},{2.5*sin(45+165)});

    \draw[given ray] (A) -- (B);
    \draw[ray] (A) -- (C);
    \node[vertex dot] at (A) {};

    \pic[draw=red!70!black, line width=0.8pt, angle radius=0.65cm,
         angle eccentricity=1.7, "$165°$", font=\small, color=red!70!black]
        {angle = B--A--C};

    \node[above right] at (1.5, 1.8) {比平角小$15°$};
    \node[below right, font=\small, color=gray] at (0.2, -0.3) {($165°$)};
  \end{scope}

\end{tikzpicture}

\vspace{1cm}
\hrule
\vspace{0.3cm}

{\small\color{gray}
\textbf{作图说明：}\\
① $45°$角：用量角器，以水平射线为基准，在上方量取$45°$，画出另一条射线。\\
② $105°$角：用量角器，以竖直射线为基准，在右侧量取$105°$，画出另一条射线。\\
③ 比平角小$15°$的角：平角$= 180°$，故所求角$= 180° - 15° = 165°$；用量角器量取即可。
}


%% ==============================
%% 第(2)题：在直线上画线段CD = 2AB
%% ==============================


\newpage

\begin{center}
  {\large\bfseries 解答：在直线上画线段 $CD$，使 $CD = 2AB$}
\end{center}

\vspace{0.8cm}

% 设定 AB 长度（单位：cm，对应 TikZ 单位）
% 图中 AB 约为斜线段，A在左下，B在右上
% 用 TikZ 计算并展示作图过程

\noindent
\begin{tikzpicture}[scale=1.1, >=Stealth]

  %% ---- 已知线段 AB（展示在左上角）----
  % A$、$B 的位置（模拟题图：B在A的右偏上方）
  \coordinate (A) at (0, 0);
  \coordinate (B) at (2.0, 0.8);   % AB 斜段，模拟题图

  % 画线段 AB
  \draw[line width=1.3pt, cyan!70!blue] (A) -- (B);
  \node[vertex dot] at (A) {};
  \node[vertex dot] at (B) {};
  \node[below left,  font=\small] at (A) {$A$};
  \node[above right, font=\small] at (B) {$B$};
  \node[above, font=\small\itshape, gray] at (1.0, 1.0) {线段 $AB$};

  %% ---- 右边的直线（水平线）----
  \draw[line width=1.2pt] (3.5, 0) -- (13.0, 0);

  %% ---- 作图步骤：用圆规量取 AB 长度两次 ----
  % AB 的长度（TikZ 计算）
  % |AB| = sqrt(2.0^2 + 0.8^2) = sqrt(4.64) ≈ 2.154 (TikZ单位)
  % 令 d = 2.154（即 AB 长），在水平线上截取

  % 取 C 为水平线上的起点
  \coordinate (C) at (4.5, 0);
  % 第一段：C 到 M（即 CM = AB）
  \coordinate (M) at ($ (C) + (2.154, 0) $);
  % 第二段：M 到 D（即 MD = AB，共 CD = 2AB）
  \coordinate (D) at ($ (M) + (2.154, 0) $);

  % 顶点点
  \node[vertex dot] at (C) {};
  \node[vertex dot] at (M) {};
  \node[vertex dot] at (D) {};

  % 标注 C$、$M$、$D
  \node[below, font=\small] at (C) {$C$};
  \node[below, font=\small\color{gray}] at (M) {（中点）};
  \node[below, font=\small] at (D) {$D$};

  % 标注线段
  \draw[line width=1.5pt, red!70!black] (C) -- (D);
  % 第一段 CM 标注
  \draw[<->, gray, line width=0.6pt] (C) -- (M)
      node[midway, above, font=\scriptsize] {$=AB$};
  % 第二段 MD 标注  
  \draw[<->, gray, line width=0.6pt] (M) -- (D)
      node[midway, above, font=\scriptsize] {$=AB$};
  % 整体 CD 标注
  \draw[<->, red!80!black, line width=0.8pt]
      ($ (C)+(0,-0.55) $) -- ($ (D)+(0,-0.55) $)
      node[midway, below, font=\small, red!70!black] {$CD = 2AB$};

  %% ---- 圆规弧线（示意作图过程）----
  % 第一段：以 C 为圆心，AB 为半径画弧，交直线于 M
  \draw[dashed, orange!80!black, line width=0.8pt]
      (C) ++(0:2.154) arc[start angle=0, end angle=18, radius=2.154];
  \draw[dashed, orange!80!black, line width=0.8pt]
      (C) ++(0:2.154) arc[start angle=0, end angle=-15, radius=2.154];
  % 虚线连接 C 到弧（示意圆规半径）
  \draw[dotted, orange!60, line width=0.7pt] (C) -- ($ (C)+(18:2.154) $);

  % 第二段：以 M 为圆心，AB 为半径画弧，交直线于 D
  \draw[dashed, orange!80!black, line width=0.8pt]
      (M) ++(0:2.154) arc[start angle=0, end angle=18, radius=2.154];
  \draw[dashed, orange!80!black, line width=0.8pt]
      (M) ++(0:2.154) arc[start angle=0, end angle=-15, radius=2.154];
  \draw[dotted, orange!60, line width=0.7pt] (M) -- ($ (M)+(18:2.154) $);

\end{tikzpicture}

\vspace{1cm}
\hrule
\vspace{0.4cm}

{\small\color{gray}
\textbf{作图说明（直尺和圆规法）：}\\
① 用圆规量取线段 $AB$ 的长度（张开圆规，使两脚分别对准 $A$、$B$）。\\
② 在直线上取一点 $C$，以 $C$ 为圆心、$AB$ 为半径画弧，弧与直线的交点记为中间点。\\
③ 再以该中间点为圆心、$AB$ 为半径画弧，弧与直线的交点即为点 $D$。\\
④ 连结 $CD$，则 $CD = 2AB$。\quad（橙色虚线弧为圆规痕迹示意）
}


%% ==============================
%% 第三题：用圆规找点B，使AB = 3厘米
%% ==============================


\newpage

\begin{center}
  {\large\bfseries 解答：用圆规在直线上找点 $B$，使 $AB = 3$ 厘米}
\end{center}

\vspace{0.8cm}

\noindent
% 此图为实际尺寸示意（1 TikZ单位 = 1cm）
\begin{tikzpicture}[scale=1.0, >=Stealth]

  %% ---- 直线（与题图一致，水平蓝线）----
  \draw[line width=1.3pt, cyan!70!blue] (-5.5, 0) -- (5.5, 0);

  %% ---- 点 A（直线中央偏左）----
  \coordinate (A) at (0, 0);
  \node[vertex dot] at (A) {};
  \node[below, font=\small] at (A) {$A$};

  %% ---- 圆规弧：以 A 为圆心，半径 3cm，画弧交直线----
  % 右侧弧（过直线上方与下方）
  \draw[dashed, orange!80!black, line width=0.9pt]
      (A) ++(0:3) arc[start angle=0, end angle=22, radius=3];
  \draw[dashed, orange!80!black, line width=0.9pt]
      (A) ++(0:3) arc[start angle=0, end angle=-18, radius=3];
  % 左侧弧
  \draw[dashed, orange!80!black, line width=0.9pt]
      (A) ++(180:3) arc[start angle=180, end angle=158, radius=3];
  \draw[dashed, orange!80!black, line width=0.9pt]
      (A) ++(180:3) arc[start angle=180, end angle=202, radius=3];

  %% ---- 辅助虚线（圆规半径示意）----
  \draw[dotted, orange!60, line width=0.7pt] (A) -- ($ (A)+(20:3) $);
  \draw[dotted, orange!60, line width=0.7pt] (A) -- ($ (A)+(160:3) $);

  %% ---- 交点 B（左右各一）----
  \coordinate (Bright) at (3, 0);
  \coordinate (Bleft)  at (-3, 0);
  \node[vertex dot, fill=red!70!black] at (Bright) {};
  \node[vertex dot, fill=red!70!black] at (Bleft)  {};
  \node[below, font=\small, red!70!black] at (Bright) {$B$};
  \node[below, font=\small, red!70!black] at (Bleft)  {$B$};

  %% ---- 尺寸标注 AB = 3厘米 ----
  % 右侧 AB
  \draw[<->, red!60!black, line width=0.7pt]
      ($ (A)+(0, 0.45) $) -- ($ (Bright)+(0, 0.45) $)
      node[midway, above, font=\small, red!70!black] {$3$ 厘米};
  % 左侧 AB
  \draw[<->, red!60!black, line width=0.7pt]
      ($ (Bleft)+(0, 0.45) $) -- ($ (A)+(0, 0.45) $)
      node[midway, above, font=\small, red!70!black] {$3$ 厘米};

\end{tikzpicture}

\vspace{1cm}
\hrule
\vspace{0.4cm}

{\small\color{gray}
\textbf{作图说明（圆规法）：}\\
① 张开圆规，使两脚距离等于 $3$ 厘米（对照直尺量取）。\\
② 将圆规针脚固定在点 $A$，用铅笔脚在直线上画弧。\\
③ 弧线与直线共有\textbf{两个交点}，均可作为点 $B$，即 $AB = 3$ 厘米。\\
（图中橙色虚弧为圆规痕迹，红点为满足条件的点 $B$）
}


%% ==============================
%% 第4题：∠1 = 50°，画∠2 = 2×∠1 = 100°
%% ==============================


\newpage

\begin{center}
  {\large\bfseries 解答：$\angle 1 = 50°$，$\angle 2 = 2 \times \angle 1 = 100°$}
\end{center}

\vspace{0.6cm}

\noindent
\begin{tikzpicture}[scale=1.1, >=Stealth]

  %%============================
  %% 左图：∠1（题目给出，测量得50°）
  %%============================
  \begin{scope}[xshift=0cm]
    \coordinate (V1) at (0, 0);       % 顶点
    \coordinate (R1) at (3.2, 0);     % 水平射线端点（向右）
    \coordinate (U1) at ({3.0*cos(50)},{3.0*sin(50)}); % 斜射线端点（50°）

    % 两条射线（青色，同题图）
    \draw[given ray] (V1) -- (R1);
    \draw[given ray] (V1) -- (U1);
    \node[vertex dot] at (V1) {};

    % 角弧 + 标注
    \pic[draw=black, line width=0.8pt, angle radius=0.75cm,
         angle eccentricity=1.55, "$1$", font=\small]
        {angle = R1--V1--U1};

    % 度数标注（红色）
    \node[red!70!black, font=\small\bfseries] at (1.5, -0.5)
        {$\angle 1 = 50°$};
  \end{scope}

  %%============================
  %% 右图：∠2 = 100°（新画）
  %%============================
  \begin{scope}[xshift=7.5cm]
    \coordinate (V2) at (0, 0);
    \coordinate (R2) at (3.2, 0);
    \coordinate (U2) at ({3.0*cos(100)},{3.0*sin(100)}); % 100°方向

    % 两条射线（青色）
    \draw[given ray] (V2) -- (R2);
    \draw[ray]       (V2) -- (U2);   % 答案射线用箭头
    \node[vertex dot] at (V2) {};

    % 角弧 + 标注
    \pic[draw=black, line width=0.8pt, angle radius=0.75cm,
         angle eccentricity=1.6, "$2$", font=\small]
        {angle = R2--V2--U2};

    % 度数标注（红色）
    \node[red!70!black, font=\small\bfseries] at (1.0, -0.5)
        {$\angle 2 = 100°$};
  \end{scope}

  %% 中间等式说明
  \node[font=\normalsize] at (5.5, 1.5)
      {$\angle 2 = 2 \times 50° = 100°$};

\end{tikzpicture}

\vspace{1cm}
\hrule
\vspace{0.4cm}

{\small\color{gray}
\textbf{解题说明：}\\
① 用量角器量出左图中 $\angle 1$ 的大小，得 $\angle 1 = 50°$。\\
② 计算：$\angle 2 = 2 \times \angle 1 = 2 \times 50° = 100°$。\\
③ 在右侧以水平射线为一边，用量角器量取 $100°$，画出另一条射线，即得 $\angle 2$。
}


%% ==============================
%% 第5题：用直尺和圆规画线段CD = 3AB
%% ==============================


\newpage

\begin{center}
  {\large\bfseries 解答：用直尺和圆规画线段 $CD$，使 $CD = 3AB$}
\end{center}

\vspace{0.6cm}

\noindent
\begin{tikzpicture}[scale=1.0, >=Stealth]

  %% ---- 已知线段 AB（左上角展示）----
  % 设 AB = 2.2（TikZ单位，约2.2cm）
  \pgfmathsetmacro{\ab}{2.2}

  \coordinate (A) at (0, 2.5);
  \coordinate (B) at (\ab, 2.5);

  \draw[line width=1.3pt, cyan!70!blue] (A) -- (B);
  \node[vertex dot] at (A) {};
  \node[vertex dot] at (B) {};
  \node[below, font=\small] at (A) {$A$};
  \node[below, font=\small] at (B) {$B$};
  \node[above, font=\small\itshape, gray] at (\ab/2, 2.7) {线段 $AB$};

  %% ---- 所画线段 CD（下方，水平）----
  % C 为起点，然后每隔AB长度打一个中间点
  \coordinate (C)  at (0,    0);
  \coordinate (M1) at (\ab,  0);          % CM1 = AB
  \coordinate (M2) at ({2*\ab}, 0);       % CM2 = 2AB
  \coordinate (D)  at ({3*\ab}, 0);       % CD  = 3AB

  % 画底线（稍长于CD，模拟直尺画出的直线）
  \draw[line width=1.0pt, gray!50] (-0.4, 0) -- ({3*\ab+0.8}, 0);

  % 红色线段CD（答案）
  \draw[line width=1.8pt, red!70!black] (C) -- (D);

  % 各点
  \node[vertex dot]                       at (C)  {};
  \node[vertex dot, fill=gray!60]         at (M1) {};
  \node[vertex dot, fill=gray!60]         at (M2) {};
  \node[vertex dot, fill=red!70!black]    at (D)  {};

  % 标注字母
  \node[below, font=\small]              at (C)  {$C$};
  \node[below, font=\small, gray]        at (M1) {};   % 中间点不标字母
  \node[below, font=\small, gray]        at (M2) {};
  \node[below, font=\small, red!70!black] at (D) {$D$};

  %% ---- 分段双箭头（各段 = AB）----
  \draw[<->, gray, line width=0.6pt]
      ($ (C)+(0, 0.45) $) -- ($ (M1)+(0, 0.45) $)
      node[midway, above, font=\scriptsize] {$=AB$};
  \draw[<->, gray, line width=0.6pt]
      ($ (M1)+(0, 0.45) $) -- ($ (M2)+(0, 0.45) $)
      node[midway, above, font=\scriptsize] {$=AB$};
  \draw[<->, gray, line width=0.6pt]
      ($ (M2)+(0, 0.45) $) -- ($ (D)+(0, 0.45) $)
      node[midway, above, font=\scriptsize] {$=AB$};

  %% ---- 总长标注 CD = 3AB ----
  \draw[<->, red!80!black, line width=0.9pt]
      ($ (C)+(0,-0.6) $) -- ($ (D)+(0,-0.6) $)
      node[midway, below, font=\small, red!70!black] {$CD = 3AB$};

  %% ---- 圆规弧示意（以 C$、$M1$、$M2 为圆心，AB为半径）----
  % 弧1：以 C 为圆心
  \draw[dashed, orange!80!black, line width=0.8pt]
      (C) ++({0:\ab}) arc[start angle=0, end angle=18,  radius=\ab];
  \draw[dashed, orange!80!black, line width=0.8pt]
      (C) ++({0:\ab}) arc[start angle=0, end angle=-15, radius=\ab];
  \draw[dotted, orange!60, line width=0.6pt] (C) -- ($ (C)+(18:\ab) $);

  % 弧2：以 M1 为圆心
  \draw[dashed, orange!80!black, line width=0.8pt]
      (M1) ++({0:\ab}) arc[start angle=0, end angle=18,  radius=\ab];
  \draw[dashed, orange!80!black, line width=0.8pt]
      (M1) ++({0:\ab}) arc[start angle=0, end angle=-15, radius=\ab];
  \draw[dotted, orange!60, line width=0.6pt] (M1) -- ($ (M1)+(18:\ab) $);

  % 弧3：以 M2 为圆心
  \draw[dashed, orange!80!black, line width=0.8pt]
      (M2) ++({0:\ab}) arc[start angle=0, end angle=18,  radius=\ab];
  \draw[dashed, orange!80!black, line width=0.8pt]
      (M2) ++({0:\ab}) arc[start angle=0, end angle=-15, radius=\ab];
  \draw[dotted, orange!60, line width=0.6pt] (M2) -- ($ (M2)+(18:\ab) $);

\end{tikzpicture}

\vspace{1cm}
\hrule
\vspace{0.4cm}

{\small\color{gray}
\textbf{作图说明（直尺和圆规法）：}\\
① 用圆规量取线段 $AB$ 的长度。\\
② 在纸上用直尺画一条射线，取一端点为 $C$。\\
③ 以 $C$ 为圆心、$AB$ 为半径画弧，交射线于中间点（第1次截取）。\\
④ 再以该点为圆心、$AB$ 为半径画弧，交射线于第二个中间点（第2次截取）。\\
⑤ 再以第二个中间点为圆心、$AB$ 为半径画弧，交射线于点 $D$（第3次截取）。\\
⑥ 线段 $CD$ 即为所求，$CD = 3AB$。\quad（橙色虚弧为圆规痕迹示意）
}


%% ==============================
%% 每两点一条直线，共几条？（2~7个点）
%% ==============================


\newpage

\begin{center}
  {\large\bfseries 解答：按照规律接着画一画}
\end{center}

\vspace{0.4cm}

% 定义常用图形命令
\newcommand{\tikzTriangle}[1]{% 空心/实心三角形
  \begin{tikzpicture}[baseline=-0.8ex]
    \draw[line width=0.8pt, #1] (0,0) -- (0.4,0) -- (0.2,0.35) -- cycle;
  \end{tikzpicture}%
}
\newcommand{\tikzCircle}[1]{% 圆形
  \begin{tikzpicture}[baseline=-0.8ex]
    \draw[line width=0.8pt, #1] (0.15,0.15) circle (0.18);
  \end{tikzpicture}%
}
\newcommand{\tikzSquare}[1]{% 正方形
  \begin{tikzpicture}[baseline=-0.8ex]
    \draw[line width=0.8pt, #1] (0,0) rectangle (0.35,0.35);
  \end{tikzpicture}%
}
\newcommand{\tikzStar}[1]{% 星形
  \begin{tikzpicture}[baseline=-0.8ex]
    \node[#1, inner sep=0pt, font=\large] at (0,0.1) {$\bigstar$};
  \end{tikzpicture}%
}
\newcommand{\tikzSolidTriangle}[1]{% 实心三角
  \begin{tikzpicture}[baseline=-0.8ex]
    \fill[#1] (0,0) -- (0.4,0) -- (0.2,0.35) -- cycle;
  \end{tikzpicture}%
}

\vspace{0.6cm}

\noindent
(1) \tikzTriangle{} \tikzCircle{} \tikzSquare{} \tikzTriangle{} \tikzCircle{} \tikzSquare{} \tikzTriangle{} \tikzCircle{} (\tikzSquare{red!70!black}) (\tikzTriangle{red!70!black}) ……

\vspace{0.8cm}

\noindent
(2) \tikzSquare{} \tikzSquare{} \tikzCircle{} \tikzTriangle{} \tikzSquare{} \tikzSquare{} \tikzCircle{} \tikzTriangle{} \tikzSquare{} (\tikzSquare{red!70!black}) \tikzCircle{} (\tikzTriangle{red!70!black}) ……

\vspace{0.8cm}

\noindent
(3) \tikzSolidTriangle{} \tikzStar{} \tikzStar{} \tikzSolidTriangle{} \tikzSolidTriangle{} \tikzSolidTriangle{} \tikzSolidTriangle{} \tikzStar{} \tikzStar{} \tikzSolidTriangle{} \tikzStar{} (\tikzStar{red!70!black}) (\tikzSolidTriangle{red!70!black}) (\tikzSolidTriangle{red!70!black}) \tikzStar{} (\tikzSolidTriangle{red!70!black}) ……

\vspace{1.2cm}
\hrule
\vspace{0.4cm}

{\small\color{gray}
\textbf{规律分析：}\\
① \textbf{第(1)题：}观察可知，图形是以 \tikzTriangle{} \tikzCircle{} \tikzSquare{} 为一组循环出现的。前面已有 \tikzTriangle{} \tikzCircle{}，按照顺序，接下来应填 \tikzSquare{} 和 \tikzTriangle{}。

\vspace{0.3cm}
② \textbf{第(2)题：}观察可知，图形是以 \tikzSquare{} \tikzSquare{} \tikzCircle{} \tikzTriangle{} 为一组循环出现的。空格前为 \tikzSquare{}，接下来的序列缺第二个 \tikzSquare{}，因此第一个括号填 \tikzSquare{}，跳过已有的 \tikzCircle{}，第二个括号填 \tikzTriangle{}。

\vspace{0.3cm}
③ \textbf{第(3)题：}该组图形由 \tikzSolidTriangle{} \tikzStar{} \tikzStar{} 和 \tikzSolidTriangle{} \tikzSolidTriangle{} \tikzSolidTriangle{} 两组交替出现。根据现有序列末尾的 \tikzSolidTriangle{} \tikzStar{}，需先补完 \tikzStar{} 构成第一组，随后补完 \tikzSolidTriangle{} \tikzSolidTriangle{} \tikzSolidTriangle{} 构成第二组。故依次填入：\tikzStar{}、\tikzSolidTriangle{}、\tikzSolidTriangle{}、\tikzSolidTriangle{}。
}


%% ==============================
%% 球类爱好项目统计
%% ==============================


\newpage

\begin{center}
  {\large\bfseries 解答：平行与垂直的练习}
\end{center}

\vspace{0.6cm}

\noindent
{\small \textbf{题目要求：} 在互相平行的两条直线下面画“//”，互相垂直的两条直线下面画“$\perp$”。}

\vspace{0.8cm}

\begin{center}
  \begin{tikzpicture}[scale=1.2, >=Stealth]
    
    %% 第一排
    % 1. 平行 (左倾)
    \begin{scope}[xshift=0cm, yshift=4cm]
      \draw[cyan!70!blue, thick] (-0.8, 1.2) -- (0.8, -0.4);
      \draw[cyan!70!blue, thick] (-0.3, 1.3) -- (1.3, -0.3);
      \node[below] at (0.25, -0.8) {$(\ \textcolor{red!70!black}{/\!/} \ )$};
    \end{scope}

    % 2. 不平行不垂直
    \begin{scope}[xshift=3.5cm, yshift=4cm]
      \draw[cyan!70!blue, thick] (-0.5, 1.3) -- (0.3, -0.3);
      \draw[cyan!70!blue, thick] (1.2, 1.5) -- (1.5, -0.4);
      \node[below] at (0.5, -0.8) {$(\ \ \ )$};
    \end{scope}

    % 3. 平行 (右倾)
    \begin{scope}[xshift=7cm, yshift=4cm]
      \draw[cyan!70!blue, thick] (-0.3, 0.2) -- (1.5, 1.8);
      \draw[cyan!70!blue, thick] (0.2, -0.3) -- (2.0, 1.3);
      \node[below] at (0.85, -0.8) {$(\ \textcolor{red!70!black}{/\!/} \ )$};
    \end{scope}

    % 4. 普通相交
    \begin{scope}[xshift=10.5cm, yshift=4cm]
      \draw[cyan!70!blue, thick] (-1.0, 1.0) -- (1.0, 0.2);
      \draw[cyan!70!blue, thick] (-0.7, -0.2) -- (0.3, 1.5);
      \node[below] at (0, -0.8) {$(\ \ \ )$};
    \end{scope}

    %% 第二排
    % 5. 普通相交 (X型)
    \begin{scope}[xshift=0cm, yshift=0cm]
      \draw[cyan!70!blue, thick] (-1.0, 1.0) -- (1.0, -0.5);
      \draw[cyan!70!blue, thick] (-1.0, -0.5) -- (1.0, 1.0);
      \node[below] at (0, -0.8) {$(\ \ \ )$};
    \end{scope}

    % 6. 垂直 (T型)
    \begin{scope}[xshift=3.5cm, yshift=0cm]
      \draw[cyan!70!blue, thick] (-1.0, 0) -- (1.0, 0.1);
      \draw[cyan!70!blue, thick] (0, 0.05) -- (0, 1.2);
      \node[below] at (0, -0.8) {$(\ \textcolor{red!70!black}{\perp} \ )$};
    \end{scope}

    % 7. 垂直 (相接)
    \begin{scope}[xshift=7cm, yshift=0cm]
      \draw[cyan!70!blue, thick] (-1.0, -0.2) -- (1.0, 0.5);
      \draw[cyan!70!blue, thick] (0.2, 0.22) -- (-0.5, 1.2);
      \node[below] at (0, -0.8) {$(\ \ \ )$};
    \end{scope}

    % 8. 平行 (陡峭)
    \begin{scope}[xshift=10.5cm, yshift=0cm]
      \draw[cyan!70!blue, thick] (-0.3, 0) -- (0.3, 2.0);
      \draw[cyan!70!blue, thick] (0.5, 0.5) -- (1.5, 2.0);
      \node[below] at (0.3, -0.8) {$(\ \ \ )$};
    \end{scope}

  \end{tikzpicture}
\end{center}

\vspace{1.0cm}
\hrule
\vspace{0.4cm}

% 逻辑说明
{\small
\textbf{判定方法说明：} \\
1. \textbf{平行：} 在同一平面内，永不相交的两条直线叫做平行线。图中第1, 3$、$8对直线满足平行关系。 \\
2. \textbf{垂直：} 当两条直线相交成直角（90$^\circ$）时，这两条直线互相垂直。图中第6, 7对直线相交成直角。 \\
3. \textbf{普通相交：} 第2, 4$、$5对直线既不平行也不垂直。
}


%% ==============================
%% 在方格纸上画平行线和垂直线
%% ==============================


\newpage

\begin{center}
  {\large\bfseries 解答：在方格纸上画平行线和垂直线}
\end{center}

\vspace{0.6cm}

\noindent
{\small \textbf{题目要求：} 在方格纸上分别画一组互相平行和互相垂直的直线。}

\vspace{0.8cm}

\begin{center}
  \begin{tikzpicture}[scale=1.0]
    % 绘制方格背景 (正方形方格)
    \draw[step=0.5cm, cyan!30, line width=0.4pt] (0,0) grid (10,4.5);
    
    % 1. 绘制一组互相平行的直线
    \draw[red!70!black, ultra thick] (1, 1) -- (4, 1);
    \draw[red!70!black, ultra thick] (1, 2) -- (4, 2);
    \node[below, font=\small, red!70!black] at (2.5, 0.8) {互相平行};

    % 2. 绘制一组互相垂直的直线
    \draw[red!70!black, ultra thick] (6, 1.5) -- (9, 1.5);
    \draw[red!70!black, ultra thick] (7.5, 0.5) -- (7.5, 3.5);
    \node[below, font=\small, red!70!black] at (7.5, 0.3) {互相垂直};
    
    % 标注直角符号
    \draw[red!70!black, thick] (7.5, 1.7) -- (7.7, 1.7) -- (7.7, 1.5);
  \end{tikzpicture}
\end{center}

\vspace{1.0cm}
\hrule
\vspace{0.4cm}

% 说明
{\small
\textbf{作作图说明：} \\
1. \textbf{平行线：} 利用方格纸的横线（或竖线）可以方便地画出平行线。图中选择了两条水平格线，它们在同一平面内且永不相交。 \\
2. \textbf{垂直线：} 利用方格纸横线与竖线互相垂直的特点，分别沿横向和纵向格线作图，相交处成90$^\circ$直角。
}


%% ==============================
%% 经过点 A 画已知直线的垂线
%% ==============================


\newpage

\begin{center}
  {\large\bfseries 解答：经过点 A 画已知直线的垂线}
\end{center}

\vspace{0.6cm}

\noindent
{\small \textbf{题目要求：} 经过点 $A$ 画已知直线的垂线。}

\vspace{0.8cm}

\begin{center}
  \begin{tikzpicture}[scale=1.2, >=Stealth]
    
    % 1. 点 A 在直线上 (斜线)
    \begin{scope}[xshift=0cm]
      % 已知直线 (斜)
      \draw[cyan!70!blue, line width=1.2pt] (-1, -1) -- (1, 1);
      % 点 A
      \node[vertex dot] (A1) at (0,0) {};
      \node[left=4pt, font=\small] at (0,0) {$A$};
      % 垂线
      \draw[red!70!black, line width=1.0pt] (-1, 1) -- (1, -1);
      % 直角符号
      \draw[red!70!black, thin] (0.15, -0.15) -- (0.3, 0) -- (0.15, 0.15);
    \end{scope}

    % 2. 点 A 在直线外 (水平线)
    \begin{scope}[xshift=5cm]
      % 已知直线 (水平)
      \draw[cyan!70!blue, line width=1.2pt] (-1.5, 0) -- (1.5, 0);
      % 点 A
      \node[vertex dot] (A2) at (0,1.5) {};
      \node[left=4pt, font=\small] at (0,1.5) {$A$};
      % 垂线
      \draw[red!70!black, line width=1.0pt] (0, 1.8) -- (0, -0.5);
      % 直角符号
      \draw[red!70!black, thin] (0, 0.2) -- (0.2, 0.2) -- (0.2, 0);
    \end{scope}

    % 3. 点 A 在直线外 (斜线)
    \begin{scope}[xshift=10cm]
      % 已知直线 (斜)
      \draw[cyan!70!blue, line width=1.2pt] (-1, -0.5) -- (1, 1);
      % 点 A
      \node[vertex dot] (A3) at (0.8, -0.4) {};
      \node[right=4pt, font=\small] at (0.8, -0.4) {$A$};
      % 垂线
      % 经过 A(0.8, -0.4) 且垂直于已知直线 (y = 0.75x + 0.25) 的垂足点是 P(0.2, 0.4)
      \draw[red!70!black, line width=1.0pt] (1.1, -0.8) -- (-0.25, 1.0);
      % 直角符号 (在此处以垂足 P(0.2, 0.4) 为锚点旋转)
      \begin{scope}[shift={(0.2, 0.4)}, rotate=36.87]
        \draw[red!70!black, thin] (0, 0.2) -- (0.2, 0.2) -- (0.2, 0);
      \end{scope}
    \end{scope}

  \end{tikzpicture}
\end{center}

\vspace{1.0cm}
\hrule
\vspace{0.4cm}

% 说明
{\small
\textbf{作图步骤说明：} \\
1. \textbf{对齐：} 用三角尺的一条直角边与已知直线重合。 \\
2. \textbf{平移：} 沿直线平移三角尺，使另一条直角边经过点 $A$。 \\
3. \textbf{作图：} 沿经过点 $A$ 的直角边画一条直线，这条直线就是已知直线的垂线。 \\
4. \textbf{标注：} 在垂足处画上直角符号“$\urcorner$”。
}


%% ==============================
%% 用画垂线的方法补画长方形和正方形
%% ==============================


\newpage

\begin{center}
  {\large\bfseries 解答：补画成长方形和正方形}
\end{center}

\vspace{0.6cm}

\noindent
{\small \textbf{题目要求：} 用画垂线的方法，将下面的图形分别补画成一个长方形和一个正方形。}

\vspace{0.8cm}

\begin{center}
  \begin{tikzpicture}[scale=1.2, >=Stealth]
    
    % 1. 补画长方形
    \begin{scope}[xshift=0cm]
      % 已知部分 (L型)
      \draw[cyan!70!blue, line width=1.2pt] (0, 1.5) -- (0, 0) -- (4, 0);
      
      % 补画部分 (红)
      \draw[red!70!black, line width=1.0pt, dashed] (4, 0) -- (4, 1.5);
      \draw[red!70!black, line width=1.0pt, dashed] (0, 1.5) -- (4, 1.5);
      
      % 直角符号 (四个角)
      \draw[cyan!70!blue, thin] (0, 0.2) -- (0.2, 0.2) -- (0.2, 0);
      \draw[red!70!black, thin] (4, 0.2) -- (3.8, 0.2) -- (3.8, 0);
      \draw[red!70!black, thin] (4, 1.3) -- (3.8, 1.3) -- (3.8, 1.5);
      \draw[red!70!black, thin] (0, 1.3) -- (0.2, 1.3) -- (0.2, 1.5);
      
      \node[below] at (2, -0.4) {补画成长方形};
    \end{scope}

    % 2. 补画正方形
    \begin{scope}[xshift=7cm]
      % 已知部分 (L型)
      \draw[cyan!70!blue, line width=1.2pt] (0, 2) -- (0, 0) -- (2, 0);
      
      % 补画部分 (红)
      \draw[red!70!black, line width=1.0pt, dashed] (2, 0) -- (2, 2);
      \draw[red!70!black, line width=1.0pt, dashed] (0, 2) -- (2, 2);
      
      % 直角符号 (四个角)
      \draw[cyan!70!blue, thin] (0, 0.2) -- (0.2, 0.2) -- (0.2, 0);
      \draw[red!70!black, thin] (2, 0.2) -- (1.8, 0.2) -- (1.8, 0);
      \draw[red!70!black, thin] (2, 1.8) -- (1.8, 1.8) -- (1.8, 2);
      \draw[red!70!black, thin] (0, 1.8) -- (0.2, 1.8) -- (0.2, 2);
      
      \node[below] at (1, -0.4) {补画成正方形};
    \end{scope}

  \end{tikzpicture}
\end{center}

\vspace{1.0cm}
\hrule
\vspace{0.4cm}

% 说明
{\small
\textbf{作图说明：} \\
1. \textbf{长方形：} 从已知线段的两个端点出发，分别画出已知线段的垂线，使垂线段的长度分别等于对应边的长度，最后连接两个垂足。 \\
2. \textbf{正方形：} 正方形是特殊的长方形，四条边都相等。由于已知两边均为2单位长度，补画的两条边也应设为2单位长度。
}


%% ==============================
%% 分别过直线外一点，画已知直线的平行线
%% ==============================


\newpage

\begin{center}
  {\large\bfseries 解答：过直线外一点画平行线}
\end{center}

\vspace{0.6cm}

\noindent
{\small \textbf{题目要求：} 分别过直线外一点，画已知直线的平行线。}

\vspace{0.8cm}

\begin{center}
  \begin{tikzpicture}[scale=1.2, >=Stealth]
    
    % 1. 水平线，点在下方
    \begin{scope}[xshift=0cm]
      % 已知直线
      \draw[cyan!70!blue, line width=1.2pt] (-1.5, 1.2) -- (1.5, 1.2);
      % 点 A
      \node[vertex dot] (A1) at (0, -0.4) {};
      \node[right=4pt, font=\small] at (0, -0.4) {$A$};
      % 平行线 (水平)
      \draw[red!70!black, line width=1.0pt] (-1.5, -0.4) -- (1.5, -0.4);
    \end{scope}

    % 2. 斜线 (中等斜率)，点在下方
    \begin{scope}[xshift=5cm]
      % 已知直线 (斜率 0.5)
      \draw[cyan!70!blue, line width=1.2pt] (-1.2, 0) -- (0.8, 1.0);
      % 点 A
      \node[vertex dot] (A2) at (0.5, -0.5) {};
      \node[right=4pt, font=\small] at (0.5, -0.5) {$A$};
      % 平行线 (斜率 0.5)
      % y = 0.5x - 0.75 => (-1.2, -1.35) 到 (1.2, -0.15)
      \draw[red!70!black, line width=1.0pt] (-1.0, -1.25) -- (1.2, -0.15);
    \end{scope}

    % 3. 斜线 (较陡斜率)，点在上方
    \begin{scope}[xshift=10cm]
      % 已知直线 (斜率 1.0)
      \draw[cyan!70!blue, line width=1.2pt] (-1, -1.5) -- (0.5, 0);
      % 点 A
      \node[vertex dot] (A3) at (-0.5, 1) {};
      \node[right=4pt, font=\small] at (-0.5, 1) {$A$};
      % 平行线 (斜率 1.0)
      % y = x + 1.5 => (-1.2, 0.3) 到 (0.3, 1.8)
      \draw[red!70!black, line width=1.0pt] (-1.2, 0.3) -- (0.3, 1.8);
    \end{scope}

  \end{tikzpicture}
\end{center}

\vspace{1.0cm}
\hrule
\vspace{0.4cm}

% 说明
{\small
\textbf{作图步骤说明：} \\
1. \textbf{对齐：} 用三角尺的一条直角边与已知直线重合。 \\
2. \textbf{靠紧：} 将直尺紧靠三角尺的另一条直角边。 \\
3. \textbf{平移：} 沿直尺平移三角尺，使原本与已知直线重合的直角边经过已知的点 $A$。 \\
4. \textbf{作图：} 沿三角尺经过点 $A$ 的这条边画一条直线，这条直线就是已知直线的平行线。
}


%% ==============================
%% 找出图形中的平行线并描红
%% ==============================


\newpage

\begin{center}
  {\large\bfseries 解答：找出图形中的平行线并描红}
\end{center}

\vspace{0.6cm}

\noindent
{\small \textbf{题目要求：} 下面的图形中，哪两条线段互相平行？用同一种颜色的笔描出来。}

\vspace{0.8cm}

\begin{center}
  \begin{tikzpicture}[scale=1.2, >=Stealth]
    
    % 1. 长方形 (旋转 20度)
    \begin{scope}[xshift=0cm, rotate=20]
      % 原始边 (青色)
      \draw[cyan!70!blue, thin] (0, 0) rectangle (3.5, 1.8);
      % 第一组平行边 (描红 - 深红)
      \draw[red!80!black, ultra thick] (0, 0) -- (3.5, 0);
      \draw[red!80!black, ultra thick] (0, 1.8) -- (3.5, 1.8);
      % 第二组平行边 (黄色)
      \draw[yellow!95!black, ultra thick] (0, 0) -- (0, 1.8);
      \draw[yellow!95!black, ultra thick] (3.5, 0) -- (3.5, 1.8);
      \node[below, rotate=20] at (1.75, -0.3) {长方形 (两组平行)};
    \end{scope}



    % 2. 正方形 (旋转 -15度)
    \begin{scope}[xshift=6cm, rotate=-15]
      % 原始边 (青色)
      \draw[cyan!70!blue, thin] (0, 0) rectangle (2.2, 2.2);
      % 第一组平行边 (描红 - 深红)
      \draw[red!80!black, ultra thick] (0, 0) -- (2.2, 0);
      \draw[red!80!black, ultra thick] (0, 2.2) -- (2.2, 2.2);
      % 第二组平行边 (黄色)
      \draw[yellow!95!black, ultra thick] (0, 0) -- (0, 2.2);
      \draw[yellow!95!black, ultra thick] (2.2, 0) -- (2.2, 2.2);
      \node[below, rotate=-15] at (1.1, -0.3) {正方形 (两组平行)};
    \end{scope}

    % 3. 梯形 (等腰梯形)
    \begin{scope}[xshift=10.5cm]
      % 原始边 (青色)
      \draw[cyan!70!blue, thin] (0, 0) -- (3, 0) -- (2.5, 1.5) -- (0.5, 1.5) -- cycle;
      % 平行边 (描红: 仅上下底)
      \draw[red!80!black, ultra thick] (0, 0) -- (3, 0);
      \draw[red!80!black, ultra thick] (0.5, 1.5) -- (2.5, 1.5);
      \node[below] at (1.5, -0.3) {梯形 (一组平行)};
    \end{scope}

  \end{tikzpicture}
\end{center}

\vspace{1.0cm}
\hrule
\vspace{0.4cm}

% 说明
{\small
\textbf{解析说明：} \\
1. \textbf{长方形和正方形：} 它们的对边分别平行。因此，长方形和正方形的四条边可以分成两组，每组包含两条互相平行的线段。图中已将这两组平行线段全部描红。 \\
2. \textbf{梯形：} 只有一组对边互相平行（上底和下底）。腰（两侧的线段）不平行。因此，图中仅描红了上下底。
}


%% ==============================
%% 画出与已知直线距离是1cm$、$2cm$、$3cm的平行线
%% ==============================


\newpage

\begin{center}
  {\large\bfseries 解答：按指定距离画平行线}
\end{center}

\vspace{0.6cm}

\noindent
{\small \textbf{题目要求：} 分别画出与已知直线距离是 1 厘米$、$2 厘米、3 厘米的平行线。}

\vspace{0.8cm}

\begin{center}
  \begin{tikzpicture}[scale=1.0, >=Stealth]
    
    % 1. 水平线 (1, 2, 3 cm)
    \begin{scope}[xshift=0cm]
      % 已知直线 (y=0)
      \draw[cyan!70!blue, line width=1.2pt] (-1.5, 0) -- (1.5, 0);
      
      % 平行线 (上方)
      \draw[red!70!black, line width=0.8pt] (-1.5, 1) -- (1.5, 1);
      \draw[red!70!black, line width=0.8pt] (-1.5, 2) -- (1.5, 2);
      \draw[red!70!black, line width=0.8pt] (-1.5, 3) -- (1.5, 3);
      
      % 距离标注 (虚线箭头)
      \draw[gray, dashed, <->] (0, 0) -- (0, 1) node[midway, right, font=\tiny] {1cm};
      \draw[gray, dashed] (0, 1) -- (0, 3);
      \draw[gray, dashed, ->] (0, 1) -- (0, 2) node[pos=1, right, font=\tiny] {2cm};
      \draw[gray, dashed, ->] (0, 2) -- (0, 3) node[pos=1, right, font=\tiny] {3cm};
      
      %\node[below] at (0, -0.2) {场景 1: 水平线};
    \end{scope}

    % 2. 斜线 (中等斜率)，1, 2, 3 cm
    \begin{scope}[xshift=5.5cm]
      % 已知直线 (斜角 30度)
      \draw[cyan!70!blue, line width=1.2pt] (-1.5, -0.866) -- (1.5, 0.866);
      
      % 平行线 (垂直位移)
      % 法向角是 30 + 90 = 120 度
      \foreach \d in {1, 2, 3} {
        \draw[red!70!black, line width=0.8pt, shift={(120:\d)}] (-1.5, -0.866) -- (1.5, 0.866);
      }
      
      % 距离标注
      \draw[gray, dashed, ->] (0,0) -- (120:1) node[midway, right=1pt, font=\tiny] {1cm};
      \draw[gray, dashed, ->] (120:1) -- (120:2) node[midway, right=1pt, font=\tiny] {2cm};
      \draw[gray, dashed, ->] (120:2) -- (120:3) node[midway, right=1pt, font=\tiny] {3cm};

      %\node[below] at (0, -1.2) {场景 2: 斜线 (30$^\circ$)};
    \end{scope}

    % 3. 斜线 (较陡斜率)，1, 2, 3 cm
    \begin{scope}[xshift=10.5cm]
      % 已知直线 (斜角 -40度)
      \draw[cyan!70!blue, line width=1.2pt] (-1.5, 1.258) -- (1.5, -1.258);
      
      % 法向角是 -40 + 90 = 50 度
      \foreach \d in {1, 2, 3} {
        \draw[red!70!black, line width=0.8pt, shift={(50:\d)}] (-1.5, 1.258) -- (1.5, -1.258);
      }
      
      % 距离标注
      \draw[gray, dashed, ->] (0,0) -- (50:1) node[midway, left=1pt, font=\tiny] {1cm};
      \draw[gray, dashed, ->] (50:1) -- (50:2) node[midway, left=1pt, font=\tiny] {2cm};
      \draw[gray, dashed, ->] (50:2) -- (50:3) node[midway, left=1pt, font=\tiny] {3cm};

      %\node[below] at (0, -1.6) {场景 3: 斜线 (-40$^\circ$)};
    \end{scope}

  \end{tikzpicture}
\end{center}

\vspace{1.0cm}
\hrule
\vspace{0.4cm}

% 说明
{\small
\textbf{作图方法说明：} \\
1. \textbf{定垂线：} 在已知直线上找一点，画出已知直线的垂线。 \\
2. \textbf{取长度：} 在垂线上分别量出 1 厘米$、$2 厘米、3 厘米的长度，标出三个点。 \\
3. \textbf{画平行：} 过这三个点，分别画出已知直线的平行线。 \\
\textit{注：平行线也可以画在已知直线的另一侧，图中仅展示了一侧的画法。}
}


%% ==============================
%% 过点 A 画垂线和平行线，并量距离
%% ==============================


\newpage

\begin{center}
  {\large\bfseries 解答：长方形和正方形拼图}
\end{center}

\vspace{0.6cm}

\noindent
{\small \textbf{题目要求：} 用一个长 6 厘米、宽 3 厘米的长方形和两个边长 3 厘米的正方形，可以拼成一个长方形，也可以拼成一个正方形。先在方格纸上画一画，再填空。}

\vspace{0.8cm}

\noindent
\textbf{1. 拼成长方形：} (可以拼成一个长 \textbf{12} 厘米、宽 \textbf{3} 厘米的长方形)

\begin{center}
  \begin{tikzpicture}[scale=0.5]
    % 网格
    \draw[cyan!30, thin] (0,0) grid (15, 5);
    
    % 拼成的长方形 (12 x 3)
    % 原始 6x3 长方形
    \fill[cyan!20] (1,1) rectangle (7,4);
    \draw[cyan!70!blue, thick] (1,1) rectangle (7,4);
    
    % 第一个 3x3 正方形
    \fill[red!20] (7,1) rectangle (10,4);
    \draw[red!70!black, thick] (7,1) rectangle (10,4);
    
    % 第二个 3x3 正方形
    \fill[yellow!20] (10,1) rectangle (13,4);
    \draw[yellow!70!black, thick] (10,1) rectangle (13,4);
    
    % 总外框 (不单独加外框，保持内部颜色区分)
    \node[below] at (7, 1) {\small 拼成长方形: $12 \times 3$};
  \end{tikzpicture}
\end{center}

\vspace{1.0cm}

\noindent
\textbf{2. 拼成正方形：} (可以拼成一个边长 \textbf{6} 厘米的正方形)

\begin{center}
  \begin{tikzpicture}[scale=0.5]
    % 网格
    \draw[cyan!30, thin] (0,0) grid (8, 8);
    
    % 拼成的正方形 (6 x 6)
    % 原始 6x3 长方形 (置于上方)
    \fill[cyan!20] (1,4) rectangle (7,7);
    \draw[cyan!70!blue, thick] (1,4) rectangle (7,7);
    
    % 第一个 3x3 正方形 (左下方)
    \fill[red!20] (1,1) rectangle (4,4);
    \draw[red!70!black, thick] (1,1) rectangle (4,4);
    
    % 第二个 3x3 正方形 (右下方)
    \fill[yellow!20] (4,1) rectangle (7,4);
    \draw[yellow!70!black, thick] (4,1) rectangle (7,4);
    
    \node[below] at (4, 1) {\small 拼成正方形: $6 \times 6$};
  \end{tikzpicture}
\end{center}

\vspace{1.0cm}
\hrule
\vspace{0.4cm}

% 说明
{\small
\textbf{解答填空：} \\
1. 可以拼成一个长 (\textbf{ 12 }) 厘米、宽 (\textbf{ 3 }) 厘米的长方形。 \\
2. 也可以拼成一个边长 (\textbf{ 6 }) 厘米的正方形。
}


%% ==============================
%% 利用直尺和圆规画图 (保留作图痕迹)
%% ==============================


\newpage

\begin{center}
  {\large\bfseries 解答：利用直尺和圆规画图 (保留作图痕迹)}
\end{center}

\vspace{0.6cm}

\begin{center}
  \begin{tikzpicture}[scale=0.8, >=Stealth]
    % (1) 画长 5cm, 宽 3cm 的长方形
    \begin{scope}[xshift=0cm]
      \coordinate (A) at (0,0);
      
      % 第一步: 画射线并作垂线
      \draw[cyan!70!blue, line width=1.0pt, ->] (-1,0) -- (6,0) node[right] {射线};
      \node[below left] at (A) {$A$};
      
      % 垂线痕迹 (以 A 为圆心截两点，再截上方)
      \draw[gray, thin] (0.8,0) arc (0:180:0.8);
      \coordinate (P1) at (-0.8,0);
      \coordinate (P2) at (0.8,0);
      \draw[gray, thin] (P1) ++(0.707, 0.707) arc (45:75:1);
      \draw[gray, thin] (P2) ++(-0.707, 0.707) arc (135:105:1);
      \draw[gray, dashed, thin] (0, -0.5) -- (0, 4) node[above] {垂线};
      
      % 第二步: 截取 B (5cm) 和 D (3cm)
      \draw[gray, thin] (5, 0.5) arc (90:45:0.7); % A 为圆心，R=5 截射线于 B
      \coordinate (B) at (5,0);
      \node[below] at (B) {$B$};
      
      \draw[gray, thin] (0.5, 3) arc (0:-35:0.7); % A 为圆心，R=3 截垂线于 D
      \coordinate (D) at (0,3);
      \node[left] at (D) {$D$};
      
      % 第三步: 以 B 为圆心 (R=3), D 为圆心 (R=5) 画弧相交于 C
      \draw[gray, thin] (D) ++(4.3, 0.3) arc (5:-5:5); % D 为圆心 R=5
      \draw[gray, thin] (B) ++(0.3, 2.3) arc (85:105:3); % B 为圆心 R=3 (修正: B 到 C 是宽，R=3)
      \coordinate (C) at (5,3);
      \node[above right] at (C) {$C$};
      
      % 第四步: 连接 BC, CD
      \draw[red!70!black, ultra thick] (A) -- (B) -- (C) -- (D) -- cycle;
      
      \node[below=1.2cm] at (2.5, 0) {\small (1) $5 \times 3$ 长方形};
    \end{scope}

    \begin{scope}[xshift=9cm]
      \coordinate (A) at (0,0);
      
      % 第一步: 画射线并作垂线
      \draw[cyan!70!blue, line width=1.0pt, ->] (-0.8,0) -- (3,0);
      \node[below left] at (A) {$A$};
      
      % 垂线痕迹
      \draw[gray, thin] (0.4,0) arc (0:180:0.4);
      \draw[gray, thin] (-0.4,0) ++(0.4, 0.6) arc (60:90:0.75);
      \draw[gray, thin] (0.4,0) ++(-0.4, 0.6) arc (120:90:0.75);
      \draw[gray, dashed, thin] (0, -0.5) -- (0, 3);
      
      % 第二步: 截取 B (2cm) 和 D (2cm)
      \draw[gray, thin] (2, 0.4) arc (90:30:0.5);
      \coordinate (B) at (2,0);
      \node[below] at (B) {$B$};
      
      \draw[gray, thin] (0.4, 2) arc (0:-30:0.5);
      \coordinate (D) at (0,2);
      \node[left] at (D) {$D$};
      
      % 第三步: 以 B, D 为圆心 (R=2) 画弧相交于 C
      \draw[gray, thin] (B) ++(0.3, 1.3) arc (80:110:2);
      \draw[gray, thin] (D) ++(1.3, 0.3) arc (10:-10:2);
      \coordinate (C) at (2,2);
      \node[above right] at (C) {$C$};
      
      % 第四步: 连接 BC, CD
      \draw[red!70!black, ultra thick] (A) -- (B) -- (C) -- (D) -- cycle;
      
      \node[below=1.2cm] at (1, 0) {\small (2) $2 \times 2$ 正方形};
    \end{scope}
  \end{tikzpicture}
\end{center}

\vspace{1.0cm}
\hrule
\vspace{0.4cm}

% 说明
{\small
\textbf{作图原理：} \\
1. \textbf{垂直：} 利用圆规在定线上截取等距两点，再以这两点为圆心，大于半径的一半为半径画弧，交点连线即为垂线。 \\
2. \textbf{度量：} 用圆规或直尺在垂线上量取指定高度，确定顶点。 \\
3. \textbf{封闭：} 连接各顶点，完成图形，并保留所有辅助圆弧痕迹。
}


\newpage

\newpage

\subsection{第 1 题：分数等值涂色}

\textbf{1.} 先在图中涂色表示已知分数的等值分数，再填空。

\vspace{0.5cm}
\begin{center}
\begin{tikzpicture}[>={Stealth[length=6pt]}, line width=0.8pt]
  \def\W{5.0}  % 横向长矩形总长
  \def\H{1.0}  % 矩形高度
  \def\gapY{-1.5} % 上下图形间隙
  \def\gapX{6.5} % 左右两侧题目水平间距
  
  % ======================= 左侧题 =======================
  \begin{scope}[shift={(0,0)}]
    % 1. 顶部矩形 (1/3)
    \def\wL{\W/3}
    \fill[gray!30] (0, 0) rectangle (\wL, \H);
    \draw (0,0) rectangle (\W, \H);
    \foreach \i in {1,2} {
      \draw[dashed, line width=0.6pt] ({\i*\wL}, 0) -- ({\i*\wL}, \H);
    }
    
    % 2. 底部矩形 (2/6)
    \def\wLb{\W/6}
    \begin{scope}[shift={(0,\gapY)}]
      \fill[green] (0, 0) rectangle (2*\wLb, \H);
      \draw (0,0) rectangle (\W, \H);
      \foreach \i in {1,...,5} {
        \draw[dashed, line width=0.6pt] ({\i*\wLb}, 0) -- ({\i*\wLb}, \H);
      }
    \end{scope}
    
    % 3. 底部公式
    \node at (\W/2, \gapY - 1.2) {\Large $\dfrac{1}{3} = \dfrac{(\;\textcolor{blue}{2}\;)}{(\;\textcolor{blue}{6}\;)}$};
  \end{scope}

  % ======================= 右侧题 =======================
  \begin{scope}[shift={(\gapX,0)}]
    % 1. 顶部矩形 (1/4)
    \def\wR{\W/4}
    \fill[gray!30] (0, 0) rectangle (\wR, \H);
    \draw (0,0) rectangle (\W, \H);
    \foreach \i in {1,2,3} {
      \draw[dashed, line width=0.6pt] ({\i*\wR}, 0) -- ({\i*\wR}, \H);
    }
    
    % 2. 底部矩形 (2/8)
    \def\wRb{\W/8}
    \begin{scope}[shift={(0,\gapY)}]
        \fill[green] (0, 0) rectangle (2*\wRb, \H);
        \draw (0,0) rectangle (\W, \H);
        \foreach \i in {1,...,7} {
          \draw[dashed, line width=0.6pt] ({\i*\wRb}, 0) -- ({\i*\wRb}, \H);
        }
    \end{scope}

    % 3. 底部公式
    \node at (\W/2, \gapY - 1.2) {\Large $\dfrac{1}{4} = \dfrac{(\;\textcolor{blue}{2}\;)}{(\;\textcolor{blue}{8}\;)}$};
  \end{scope}

\end{tikzpicture}
\end{center}

\vspace{0.5cm}
\textbf{绘图原理：}
\begin{enumerate}
    \item \textbf{布局参数：} 定义基础矩形宽 \texttt{W=5.0}、高 \texttt{H=1.0}。使用 \texttt{scope} 和 \texttt{shift} 参数将左右两题、上下两图隔开。
    \item \textbf{等值涂色：} 左侧分别按 3 等分和 6 等分规定网格尺寸，对应涂色 1 格（\texttt{gray!30}）与 2 格（\texttt{green}）。右侧同理按照 4 等分和 8 等分处理，完成 1 格与 2 格的面积图形映射。
    \item \textbf{图层管理：} 先调用 \texttt{\textbackslash fill} 填涂面积底层，再使用 \texttt{\textbackslash draw} 覆盖 0.8pt 实线外围与 0.6pt 内部虚线，防止线条被遮挡。
\end{enumerate}

\vspace{1.5cm}

\newpage

\subsection{第 2 题：分数等值涂色（多边形与正方形）}

\textbf{2.} 先在图中涂色表示已知分数的等值分数，再填空。

\begin{center}
\begin{tikzpicture}[>={Stealth[length=6pt]}, line width=0.8pt]
  % ==== 1. 第一组（六边形 1/3 = 2/6） ====
  % 1.1 第一个六边形 (1/3)
  \begin{scope}[shift={(0,0)}]
    % 涂色部分 (左侧菱形)
    \fill[gray!30] (180:1.3) -- (120:1.3) -- (0,0) -- (240:1.3) -- cycle;
    % 外边框
    \draw (0:1.3) \foreach \a in {60, 120, 180, 240, 300} { -- (\a:1.3) } -- cycle;
    % 内部虚线
    \draw[dashed, line width=0.6pt] (0,0) -- (0:1.3);
    \draw[dashed, line width=0.6pt] (0,0) -- (120:1.3);
    \draw[dashed, line width=0.6pt] (0,0) -- (240:1.3);
    % 下方分数
    \node[font=\Large] at (0, -2.5) {$\dfrac{1}{3}$};
  \end{scope}

  % 等号
  \node[font=\Large] at (2.0, -2.5) {$=$};

  % 1.2 第二个六边形 (2/6)
  \begin{scope}[shift={(4.0,0)}]
    % 涂色部分 (左侧两个三角形，对应 2/6，颜色为红）
    \fill[red] (180:1.3) -- (120:1.3) -- (0,0) -- (240:1.3) -- cycle;
    % 外边框
    \draw (0:1.3) \foreach \a in {60, 120, 180, 240, 300} { -- (\a:1.3) } -- cycle;
    % 内部虚线 (对角线连线，6等分)
    \draw[dashed, line width=0.6pt] (0:1.3) -- (180:1.3);
    \draw[dashed, line width=0.6pt] (60:1.3) -- (240:1.3);
    \draw[dashed, line width=0.6pt] (120:1.3) -- (300:1.3);
    
    % 下方分数
    \node[font=\Large] at (0, -2.5) {$\dfrac{(\;\textcolor{blue}{2}\;)}{(\;\textcolor{blue}{6}\;)}$};
  \end{scope}

  % ==== 2. 第二组（正方形 1/4 = 4/16） ====
  % 定义正方形的一半边长为 1.15
  \def\myL{1.15}
  % 距离第二六边形有一定偏移
  \begin{scope}[shift={(8.5,0)}]
    % 涂色部分 (顶部 1/4)
    \fill[gray!30] (-\myL, \myL/2) rectangle (\myL, \myL);
    % 外边框
    \draw (-\myL, -\myL) rectangle (\myL, \myL);
    % 内部虚线 (横向3条分割成4等分)
    \draw[dashed, line width=0.6pt] (-\myL, \myL/2) -- (\myL, \myL/2);
    \draw[dashed, line width=0.6pt] (-\myL, 0) -- (\myL, 0);
    \draw[dashed, line width=0.6pt] (-\myL, -\myL/2) -- (\myL, -\myL/2);
    
    % 下方分数
    \node[font=\Large] at (0, -2.5) {$\dfrac{1}{4}$};
  \end{scope}

  % 等号
  \node[font=\Large] at (10.6, -2.5) {$=$};

  % 2.2 第二个正方形 (4/16)
  \begin{scope}[shift={(12.7,0)}]
    % 涂色部分 (顶部第一排四个小正方形，即顶部 1/4)
    \fill[red] (-\myL, \myL/2) rectangle (\myL, \myL);
    % 外边框
    \draw (-\myL, -\myL) rectangle (\myL, \myL);
    % 内部虚线 (4x4网格，横竖各3条)
    % 横向
    \draw[dashed, line width=0.6pt] (-\myL, \myL/2) -- (\myL, \myL/2);
    \draw[dashed, line width=0.6pt] (-\myL, 0) -- (\myL, 0);
    \draw[dashed, line width=0.6pt] (-\myL, -\myL/2) -- (\myL, -\myL/2);
    % 纵向
    \draw[dashed, line width=0.6pt] (-\myL/2, -\myL) -- (-\myL/2, \myL);
    \draw[dashed, line width=0.6pt] (0, -\myL) -- (0, \myL);
    \draw[dashed, line width=0.6pt] (\myL/2, -\myL) -- (\myL/2, \myL);
    
    % 下方分数
    \node[font=\Large] at (0, -2.5) {$\dfrac{(\;\textcolor{blue}{4}\;)}{(\;\textcolor{blue}{16}\;)}$};
  \end{scope}
\end{tikzpicture}
\end{center}

\vspace{0.5cm}
\textbf{绘图原理：}
\begin{enumerate}
    \item \textbf{布局设置：} 设定四组图形的中心点横坐标，使用 \texttt{scope} 的 \texttt{shift} 将图形纵向居中并平铺。文字公式统一置于纵向位移 -2.5 处对齐。
    \item \textbf{六边形绘制：} 采用极坐标生成正六边形，设定起始角度为 $0^\circ$，步进 $60^\circ$。
    \item \textbf{等宽虚线划分：} 六边形内部连接对角线实现 6 等分；正方形内部通过横纵各 3 条虚线完成 $4\times 4$ 网格拆分。
    \item \textbf{分块填色：} 左侧六边形涂满占据 $1/3$ 面积的独立菱形块。右侧正方形取半边长 \texttt{L=1.15}，分别涂顶部 1/4 区域，完成 $1/4 = 4/16$ 的面积比照对照。
\end{enumerate}

\vspace{1.5cm}

\newpage

\subsection{第 2 题（补充）：长方形网格分数比例}

\textbf{2.} 把下面长方形的 $\dfrac{4}{8}$ 涂红色，$\dfrac{3}{8}$ 涂蓝色。

\vspace{0.5cm}
\begin{center}
\begin{tikzpicture}[>={Stealth[length=6pt]}, line width=0.8pt]
  % 定义整体长方形的宽和高 (比例 5.7:1.5)
  \def\W{5.7}
  \def\H{1.5}
  \def\cw{\W/4} % 单格列宽
  \def\rh{\H/2} % 单格行高
  
  % 涂红色区域 (前两列，共 4 格)
  \fill[red] (0,0) rectangle (2*\cw, \H);
  % 涂蓝色区域 (第 3 列 2 格 + 第 4 列上半 1 格，共 3 格)
  \fill[cyan] (2*\cw, \rh) rectangle (\W, \H);
  \fill[cyan] (2*\cw, 0)  rectangle (3*\cw, \rh);
  
  % 外围大实线方框
  \draw (0,0) rectangle (\W, \H);
  
  % 内部虚线网格
  % 1 条横向虚线
  \draw[dashed, line width=0.6pt] (0, \rh) -- (\W, \rh);
  % 3 条竖向虚线
  \foreach \x in {1,2,3} {
    \draw[dashed, line width=0.6pt] ({\x*\cw}, 0) -- ({\x*\cw}, \H);
  }
\end{tikzpicture}

\vspace{0.4cm}
{\Large $(\;\textcolor{blue}{4}\;)$ 个 $\dfrac{1}{8}$ 是 $\dfrac{4}{8}$，$(\;\textcolor{blue}{3}\;)$ 个 $\dfrac{1}{8}$ 是 $\dfrac{3}{8}$。}
\end{center}

\vspace{0.5cm}
\textbf{绘图原理：}
\begin{enumerate}
    \item \textbf{矩形比例：} 按照 5.7:1.5 比例构建原矩形，划分横向 4 列 \texttt{\textbackslash cw=\textbackslash W/4} 与纵向 2 行 \texttt{\textbackslash rh=\textbackslash H/2}，形成总共 8 等分的网格结构。
    \item \textbf{分区着色：} 前 4 份（前 2 列）运用 \texttt{\textbackslash fill[red]} 填充为红。其余 3 份划分为第 3、4 列上半以及第 3 列下侧的色块进行 \texttt{\textbackslash fill[cyan]} 着色，实现 $4/8$ 与 $3/8$ 涂填对照。
    \item \textbf{叠放次序：} 预先声明 \texttt{\textbackslash fill} 填充块，其后再添加深色矩形框架边界以及水平或垂直的 \texttt{\textbackslash foreach} 拆线虚设层。
\end{enumerate}

\vspace{1.5cm}

\newpage

\subsection{第 3 题：分数加减与涂色}

\textbf{3.} 把下面的长方形的 $\dfrac{4}{7}$ 涂红色，$\dfrac{2}{7}$ 涂绿色。绿色部分占
这个长方形的 $\dfrac{(\;\textcolor{blue}{2}\;)}{(\;\textcolor{blue}{7}\;)}$，红色部分比绿色部分多占这
个长方形的 $\dfrac{(\;\textcolor{blue}{2}\;)}{(\;\textcolor{blue}{7}\;)}$。

\vspace{0.5cm}
\begin{center}
\begin{tikzpicture}[>={Stealth[length=6pt]}, line width=0.8pt]
  \def\W{7.0}
  \def\H{1.5}
  \def\cw{\W/7}
  
  % 涂色区域 (底层)
  \fill[red] (0,0) rectangle (4*\cw, \H);
  \fill[green] (4*\cw, 0) rectangle (6*\cw, \H);
  
  % 外边框与内部虚线
  \draw (0,0) rectangle (\W, \H);
  \foreach \x in {1,...,6} {
    \draw[dashed, line width=0.6pt] ({\x*\cw}, 0) -- ({\x*\cw}, \H);
  }
\end{tikzpicture}
\end{center}

\vspace{0.5cm}
\textbf{绘图原理：}
\begin{enumerate}
    \item \textbf{图形参数：} 定义宽 \texttt{W=7.0}、高 \texttt{H=1.5}，并将列宽常数规限为 \texttt{W/7}，恰好将图形均匀划分为 7 列。
    \item \textbf{定距涂色：} 对 $x \in (0, 4\text{cw})$ 区间执行 \texttt{red} 涂色表示 $4/7$；针对 $x \in (4\text{cw}, 6\text{cw})$ 区间则执行 \texttt{green} 填涂表示 $2/7$；余下一格保持空白态。
    \item \textbf{图层排布：} 底色区块着色后，利用全局 \texttt{rectangle} 收束外部粗框，并借助 1 到 6 的 \texttt{\textbackslash foreach} 生成组垂直虚隔离网格。
\end{enumerate}

\vspace{1.5cm}

\newpage

\subsection{第 4 题：方格纸等效分数}

\textbf{4.} 在下面的图上涂色表示 $\dfrac{4}{12}$。涂色部分还可以表示
这张方格纸的 $\dfrac{(\;\textcolor{blue}{2}\;)}{(\;\textcolor{blue}{6}\;)}$ 和 $\dfrac{(\;\textcolor{blue}{1}\;)}{(\;\textcolor{blue}{3}\;)}$。

\vspace{0.5cm}
\begin{center}
\begin{tikzpicture}[>={Stealth[length=6pt]}, line width=0.8pt]
  \def\w{1.0} % 单个格子的边长设定为 1cm
  
  % 1. 涂色部分 (左侧两列共 4 个格子，对应 4/12)
  \fill[red] (0,0) rectangle (2*\w, 2*\w);
  
  % 2. 外部实线大边框 (6列 x 2行)
  \draw (0,0) rectangle (6*\w, 2*\w);
  
  % 3. 内部虚线网格
  % 横向中线
  \draw[dashed, line width=0.6pt] (0, \w) -- (6*\w, \w);
  % 5条纵向分界线
  \foreach \x in {1,2,3,4,5} {
    \draw[dashed, line width=0.6pt] (\x*\w, 0) -- (\x*\w, 2*\w);
  }
\end{tikzpicture}
\end{center}

\vspace{0.5cm}
\textbf{绘图原理：}
\begin{enumerate}
    \item \textbf{文本分离：} 题干及公式填空项脱离图片主体放置于 LaTeX 正文进行自然段落排版。作图环境独立在下方，避免多重嵌套节点的文本断层。
    \item \textbf{构建网格：} 设单网格宽为 \texttt{w=1.0}，设定坐标点 \texttt{(6\textbackslash w, 2\textbackslash w)} 以创造含 12 个格子的矩阵实框。横向加中心贯通线，纵向加 5 条分隔虚线完成格体成型。
    \item \textbf{色域核算：} 给定表达体量 $4/12$，以 \texttt{\textbackslash fill[red]} 在画布坐标起点对 $2\times2$ 共 4 单元实施覆填。它从水平长比端对应展示 $1/3$，向列占位比测视契合等效表示面积 $2/6$。
\end{enumerate}

\vspace{1.5cm}

\newpage

\subsection{第 4 题（综合）：比较分数大小与等值划分}

\textbf{4.} 乐乐和小林用同样大的红纸做手工。乐乐做花用去
红纸的 $\dfrac{1}{4}$，小林做小旗用去红纸的 $\dfrac{1}{5}$。谁用去的红纸
多一些？想一想，涂一涂，填一填。

\vspace{0.5cm}
\begin{center}
\begin{tikzpicture}[>={Stealth[length=6pt]}, line width=0.8pt]
  \def\W{2.5}
  \def\H{2.0}
  
  % ================= 行 1 =================
  % 1. 左侧(乐乐 1/4)
  \begin{scope}[shift={(0,0)}]
    \fill[gray!30] (0, 0.75*\H) rectangle (\W, \H); % 顶部 1/4 涂灰
    \draw (0,0) rectangle (\W, \H);
    \draw[dashed, line width=0.6pt] (0, 0.25*\H) -- (\W, 0.25*\H);
    \draw[dashed, line width=0.6pt] (0, 0.50*\H) -- (\W, 0.50*\H);
    \draw[dashed, line width=0.6pt] (0, 0.75*\H) -- (\W, 0.75*\H);
  \end{scope}

  % 箭头
  \draw[->, line width=1.5pt] (\W+0.4, 0.5*\H) -- (\W+1.0, 0.5*\H);

  % 2. 结果图 (20宫格，5/20)
  \begin{scope}[shift={(\W+1.4, 0)}]
    \fill[cyan] (0, 0.75*\H) rectangle (\W, \H); % 顶部一行涂蓝
    \draw (0,0) rectangle (\W, \H);
    % 划线
    \foreach \y in {0.25, 0.50, 0.75} {
      \draw[dashed, line width=0.6pt] (0, \y*\H) -- (\W, \y*\H);
    }
    \foreach \x in {0.2, 0.4, 0.6, 0.8} {
      \draw[dashed, line width=0.6pt] (\x*\W, 0) -- (\x*\W, \H);
    }
  \end{scope}

  % 3. 公式
  \node at (2*\W + 3, 0.5*\H) {\Large $\dfrac{1}{4} = \dfrac{(\;\textcolor{blue}{5}\;)}{20}$};

  % ================= 行 2 =================
  \def\Y{-2.8}
  % 1. 左侧(小林 1/5)
  \begin{scope}[shift={(0,\Y)}]
    \fill[gray!30] (0, 0) rectangle (0.2*\W, \H); % 左侧 1/5 涂灰
    \draw (0,0) rectangle (\W, \H);
    \foreach \x in {0.2, 0.4, 0.6, 0.8} {
      \draw[dashed, line width=0.6pt] (\x*\W, 0) -- (\x*\W, \H);
    }
  \end{scope}

  % 箭头
  \draw[->, line width=1.5pt] (\W+0.4, \Y+0.5*\H) -- (\W+1.0, \Y+0.5*\H);

  % 2. 结果图 (20宫格，4/20)
  \begin{scope}[shift={(\W+1.4, \Y)}]
    \fill[cyan] (0, 0) rectangle (0.2*\W, \H); % 左起第一列涂蓝
    \draw (0,0) rectangle (\W, \H);
    % 划线
    \foreach \y in {0.25, 0.50, 0.75} {
      \draw[dashed, line width=0.6pt] (0, \y*\H) -- (\W, \y*\H);
    }
    \foreach \x in {0.2, 0.4, 0.6, 0.8} {
      \draw[dashed, line width=0.6pt] (\x*\W, 0) -- (\x*\W, \H);
    }
  \end{scope}

  % 3. 公式
  \node at (2*\W + 3, \Y+0.5*\H) {\Large $\dfrac{1}{5} = \dfrac{(\;\textcolor{blue}{4}\;)}{20}$};

\end{tikzpicture}

\vspace{0.8cm}
{\Large 因为 $\dfrac{(\;\textcolor{blue}{5}\;)}{20} \;\tikz[baseline=-0.8ex] \node[circle, draw, inner sep=1pt, minimum size=1.4em, font=\large] {\textcolor{blue}{$>$}};\; \dfrac{(\;\textcolor{blue}{4}\;)}{20}$，所以 $\dfrac{1}{4} \;\tikz[baseline=-0.8ex] \node[circle, draw, inner sep=1pt, minimum size=1.4em, font=\large] {\textcolor{blue}{$>$}};\; \dfrac{1}{5}$。}
\end{center}

\vspace{0.5cm}
\textbf{绘图原理：}
\begin{enumerate}
    \item \textbf{总框网分定型：} 定义外长方形宽为 \texttt{W=2.5} 且 \texttt{H=2.0}，对应使得图形做五分与四分时，其内向正方格均保持统一底径标尺 $0.5$。
    \item \textbf{上方映射（代表 $1/4$）：} 左侧实行 4 等分数标，并在高度 $0.75\text{H}$ 到 $\text{H}$ 范围内添设底层灰标记。右侧延伸使用相同的区域染注为天蓝色 \texttt{cyan}，后设纵向 5 等分虚线阵扩成 20 分率格子网络体系。
    \item \textbf{下方映射（代表 $1/5$）：} 将 \texttt{y} 方向下移常数位 \texttt{Y=-2.8}。将最左列纵带漆作灰色示意 $1/5$；右侧同样切组为 20 分阵图，把同位置等分面积赋予统一比对新色。
    \item \textbf{符号化圆引节点：} 对文中标记圈符运用 \texttt{\textbackslash tikz[baseline=-0.8ex]} 组件模块化整合操作，强行维持内置符号相对于行底线的等高位置挂载。
\end{enumerate}

\vspace{1.5cm}

\newpage

\subsection{第 5 题：长条形方格等效分数}

\textbf{5.} 在下图中分一分，涂一涂，写出 $\dfrac{1}{5}$ 的两个等值分数。

\vspace{0.5cm}
\begin{center}
\begin{tikzpicture}[>={Stealth[length=6pt]}, line width=0.8pt]
  \def\W{10}   % 总宽度 10cm
  \def\H{0.8}  % 高度 0.8cm
  \def\gap{0.3} % 间距 0.3cm
  
  % ==== Strip 3 (底部，15等分) ====
  \begin{scope}[shift={(0,0)}]
    % 涂色区域 (占 1/5 长度)
    \fill[red] (0,0) rectangle (\W/5, \H);
    \draw (0,0) rectangle (\W, \H);
    % 实线 (将每一大块均分为3份)
    \foreach \b in {0,...,4} {
      \draw[line width=0.6pt] ({\b*\W/5 + \W/15}, 0) -- ({\b*\W/5 + \W/15}, \H);
      \draw[line width=0.6pt] ({\b*\W/5 + 2*\W/15}, 0) -- ({\b*\W/5 + 2*\W/15}, \H);
    }
    % 虚线 (原 5等分 线)
    \foreach \i in {1,...,4} {
      \draw[dashed, line width=0.6pt] ({\i*\W/5}, 0) -- ({\i*\W/5}, \H);
    }
  \end{scope}

  % ==== Strip 2 (中部，10等分) ====
  \begin{scope}[shift={(0,\H+\gap)}]
    \fill[red] (0,0) rectangle (\W/5, \H);
    \draw (0,0) rectangle (\W, \H);
    % 实线 (将每一大块均分为2份)
    \foreach \b in {0,...,4} {
      \draw[line width=0.6pt] ({\b*\W/5 + \W/10}, 0) -- ({\b*\W/5 + \W/10}, \H);
    }
    % 虚线 (原 5等分 线)
    \foreach \i in {1,...,4} {
      \draw[dashed, line width=0.6pt] ({\i*\W/5}, 0) -- ({\i*\W/5}, \H);
    }
  \end{scope}

  % ==== Strip 1 (顶部，原 5等分) ====
  \begin{scope}[shift={(0,2*\H+2*\gap)}]
    \fill[gray!30] (0,0) rectangle (\W/5, \H);
    \draw (0,0) rectangle (\W, \H);
    % 虚线
    \foreach \i in {1,...,4} {
      \draw[dashed, line width=0.6pt] ({\i*\W/5}, 0) -- ({\i*\W/5}, \H);
    }
  \end{scope}
\end{tikzpicture}

\vspace{0.4cm}
{\Large $\dfrac{1}{5} = \dfrac{(\;\textcolor{blue}{2}\;)}{(\;\textcolor{blue}{10}\;)} = \dfrac{(\;\textcolor{blue}{3}\;)}{(\;\textcolor{blue}{15}\;)}$}
\end{center}

\vspace{0.5cm}
\textbf{绘图原理：}
\begin{enumerate}
    \item \textbf{基本堆叠：} 三项长条通过 \texttt{shift} 将每列向上移位 \texttt{\textbackslash H + \textbackslash gap} 完成排列。尺寸统一为宽 $10\text{cm}$、高 $0.8\text{cm}$。
    \item \textbf{基础切分：} 三大矩形块的外部框线均沿 \texttt{W/5} 跨距布设等分的段落虚线。
    \item \textbf{内部重划：} 二号图附加绘制每组 $W/5$ 段内部的中分平分线；三号图附加绘制每组内部处于 $W/15$ 及 $2W/15$ 位置的实线等分截断。
    \item \textbf{赋彩分配：} 第一栏以浅灰色标记 $1/5$ 的方框；第二、三栏基于等值对应将在其同空间的起手 $W/5$ 作满覆红色预处理以示意比值均等。
\end{enumerate}

\vspace{2.5cm}
\begin{center}
\textbf{\huge 尺规作图类}
\end{center}
\vspace{1cm}
